% Options for packages loaded elsewhere
\PassOptionsToPackage{unicode}{hyperref}
\PassOptionsToPackage{hyphens}{url}
\documentclass[
]{article}
\usepackage{xcolor}
\usepackage[margin=1in]{geometry}
\usepackage{amsmath,amssymb}
\setcounter{secnumdepth}{-\maxdimen} % remove section numbering
\usepackage{iftex}
\ifPDFTeX
  \usepackage[T1]{fontenc}
  \usepackage[utf8]{inputenc}
  \usepackage{textcomp} % provide euro and other symbols
\else % if luatex or xetex
  \usepackage{unicode-math} % this also loads fontspec
  \defaultfontfeatures{Scale=MatchLowercase}
  \defaultfontfeatures[\rmfamily]{Ligatures=TeX,Scale=1}
\fi
\usepackage{lmodern}
\ifPDFTeX\else
  % xetex/luatex font selection
\fi
% Use upquote if available, for straight quotes in verbatim environments
\IfFileExists{upquote.sty}{\usepackage{upquote}}{}
\IfFileExists{microtype.sty}{% use microtype if available
  \usepackage[]{microtype}
  \UseMicrotypeSet[protrusion]{basicmath} % disable protrusion for tt fonts
}{}
\makeatletter
\@ifundefined{KOMAClassName}{% if non-KOMA class
  \IfFileExists{parskip.sty}{%
    \usepackage{parskip}
  }{% else
    \setlength{\parindent}{0pt}
    \setlength{\parskip}{6pt plus 2pt minus 1pt}}
}{% if KOMA class
  \KOMAoptions{parskip=half}}
\makeatother
\usepackage{color}
\usepackage{fancyvrb}
\newcommand{\VerbBar}{|}
\newcommand{\VERB}{\Verb[commandchars=\\\{\}]}
\DefineVerbatimEnvironment{Highlighting}{Verbatim}{commandchars=\\\{\}}
% Add ',fontsize=\small' for more characters per line
\usepackage{framed}
\definecolor{shadecolor}{RGB}{248,248,248}
\newenvironment{Shaded}{\begin{snugshade}}{\end{snugshade}}
\newcommand{\AlertTok}[1]{\textcolor[rgb]{0.94,0.16,0.16}{#1}}
\newcommand{\AnnotationTok}[1]{\textcolor[rgb]{0.56,0.35,0.01}{\textbf{\textit{#1}}}}
\newcommand{\AttributeTok}[1]{\textcolor[rgb]{0.13,0.29,0.53}{#1}}
\newcommand{\BaseNTok}[1]{\textcolor[rgb]{0.00,0.00,0.81}{#1}}
\newcommand{\BuiltInTok}[1]{#1}
\newcommand{\CharTok}[1]{\textcolor[rgb]{0.31,0.60,0.02}{#1}}
\newcommand{\CommentTok}[1]{\textcolor[rgb]{0.56,0.35,0.01}{\textit{#1}}}
\newcommand{\CommentVarTok}[1]{\textcolor[rgb]{0.56,0.35,0.01}{\textbf{\textit{#1}}}}
\newcommand{\ConstantTok}[1]{\textcolor[rgb]{0.56,0.35,0.01}{#1}}
\newcommand{\ControlFlowTok}[1]{\textcolor[rgb]{0.13,0.29,0.53}{\textbf{#1}}}
\newcommand{\DataTypeTok}[1]{\textcolor[rgb]{0.13,0.29,0.53}{#1}}
\newcommand{\DecValTok}[1]{\textcolor[rgb]{0.00,0.00,0.81}{#1}}
\newcommand{\DocumentationTok}[1]{\textcolor[rgb]{0.56,0.35,0.01}{\textbf{\textit{#1}}}}
\newcommand{\ErrorTok}[1]{\textcolor[rgb]{0.64,0.00,0.00}{\textbf{#1}}}
\newcommand{\ExtensionTok}[1]{#1}
\newcommand{\FloatTok}[1]{\textcolor[rgb]{0.00,0.00,0.81}{#1}}
\newcommand{\FunctionTok}[1]{\textcolor[rgb]{0.13,0.29,0.53}{\textbf{#1}}}
\newcommand{\ImportTok}[1]{#1}
\newcommand{\InformationTok}[1]{\textcolor[rgb]{0.56,0.35,0.01}{\textbf{\textit{#1}}}}
\newcommand{\KeywordTok}[1]{\textcolor[rgb]{0.13,0.29,0.53}{\textbf{#1}}}
\newcommand{\NormalTok}[1]{#1}
\newcommand{\OperatorTok}[1]{\textcolor[rgb]{0.81,0.36,0.00}{\textbf{#1}}}
\newcommand{\OtherTok}[1]{\textcolor[rgb]{0.56,0.35,0.01}{#1}}
\newcommand{\PreprocessorTok}[1]{\textcolor[rgb]{0.56,0.35,0.01}{\textit{#1}}}
\newcommand{\RegionMarkerTok}[1]{#1}
\newcommand{\SpecialCharTok}[1]{\textcolor[rgb]{0.81,0.36,0.00}{\textbf{#1}}}
\newcommand{\SpecialStringTok}[1]{\textcolor[rgb]{0.31,0.60,0.02}{#1}}
\newcommand{\StringTok}[1]{\textcolor[rgb]{0.31,0.60,0.02}{#1}}
\newcommand{\VariableTok}[1]{\textcolor[rgb]{0.00,0.00,0.00}{#1}}
\newcommand{\VerbatimStringTok}[1]{\textcolor[rgb]{0.31,0.60,0.02}{#1}}
\newcommand{\WarningTok}[1]{\textcolor[rgb]{0.56,0.35,0.01}{\textbf{\textit{#1}}}}
\usepackage{graphicx}
\makeatletter
\newsavebox\pandoc@box
\newcommand*\pandocbounded[1]{% scales image to fit in text height/width
  \sbox\pandoc@box{#1}%
  \Gscale@div\@tempa{\textheight}{\dimexpr\ht\pandoc@box+\dp\pandoc@box\relax}%
  \Gscale@div\@tempb{\linewidth}{\wd\pandoc@box}%
  \ifdim\@tempb\p@<\@tempa\p@\let\@tempa\@tempb\fi% select the smaller of both
  \ifdim\@tempa\p@<\p@\scalebox{\@tempa}{\usebox\pandoc@box}%
  \else\usebox{\pandoc@box}%
  \fi%
}
% Set default figure placement to htbp
\def\fps@figure{htbp}
\makeatother
\setlength{\emergencystretch}{3em} % prevent overfull lines
\providecommand{\tightlist}{%
  \setlength{\itemsep}{0pt}\setlength{\parskip}{0pt}}
\usepackage{bookmark}
\IfFileExists{xurl.sty}{\usepackage{xurl}}{} % add URL line breaks if available
\urlstyle{same}
\hypersetup{
  pdftitle={Predykcja hitów muzycznych na podstawie danych Spotify},
  hidelinks,
  pdfcreator={LaTeX via pandoc}}

\title{Predykcja hitów muzycznych na podstawie danych Spotify}
\author{}
\date{\vspace{-2.5em}2026-05-18}

\begin{document}
\maketitle

\subsection{Wybór zbioru danych}\label{wybuxf3r-zbioru-danych}

Do realizacji projektu wybrano zbiór danych o nazwie \textbf{``Spotify
Tracks Dataset''} (dostępny m.in. na platformie Kaggle). Zbiór ten
zawiera wyczerpujące metadane dotyczące tysięcy utworów muzycznych
dostępnych na platformie Spotify. Obejmuje on zróżnicowane atrybuty, w
tym parametry akustyczne wyliczane przez algorytmy Spotify (np.
taneczność, energia, akustyczność), informacje o popularności,
przynależność do gatunków muzycznych oraz podstawowe dane
identyfikacyjne (tytuł, artysta, rok wydania). Wybór tego zbioru
podyktowany jest jego bogactwem w informacje o naturze zarówno
numerycznej, jak i kategorycznej, co pozwala na zastosowanie szerokiego
spektrum technik eksploracji danych.

\begin{enumerate}
\def\labelenumi{\arabic{enumi}.}
\tightlist
\item
  Zmienne objaśniane (Target):
\end{enumerate}

\begin{itemize}
\item
  popularity (numeryczna): Skala od 0 do 100 wyliczana przez algorytmy
  Spotify, oceniająca aktualną popularność utworu (w dużej mierze oparta
  na całkowitej liczbie odtworzeń oraz tym, jak niedawno miały one
  miejsce).
\item
  is\_hit (kategoryczna/binarna): Zmienna docelowa utworzona sztucznie
  na potrzeby projektu. Przyjmuje wartość 1 (Hit), jeśli popularity
  wynosi 75 lub więcej, oraz 0 (Brak hitu) w przeciwnym wypadku.
\end{itemize}

\begin{enumerate}
\def\labelenumi{\arabic{enumi}.}
\setcounter{enumi}{1}
\tightlist
\item
  Cechy psychoakustyczne (skala 0.0 - 1.0):
\end{enumerate}

\begin{itemize}
\item
  danceability: Określa, jak bardzo utwór nadaje się do tańca na
  podstawie kombinacji elementów muzycznych, takich jak tempo,
  stabilność rytmu i siła bitu.
\item
  energy: Stanowi miarę intensywności i aktywności utworu. Utwory o
  wysokiej energii są odbierane jako szybkie, głośne i hałaśliwe (np.
  death metal ma wysoką energię).
\item
  acousticness: Wskaźnik określający pewność algorytmu, że utwór ma
  charakter czysto akustyczny (nieelektroniczny).
\item
  instrumentalness: Przewiduje, czy utwór nie zawiera partii wokalnych
  (słów). Wartości bliskie 1.0 oznaczają utwory czysto instrumentalne.
\item
  liveness: Określa prawdopodobieństwo obecności publiczności podczas
  nagrania. Wyższe wartości wskazują na wykonanie na żywo (live).
\item
  speechiness: Wykrywa obecność mówionych słów w nagraniu. Wartości
  powyżej 0.66 opisują utwory składające się niemal wyłącznie ze słów
  (np. podcasty, stand-up), natomiast wartości poniżej 0.33 to typowa
  muzyka z wokalem.
\item
  valence: Opisuje muzyczną ``pozytywność'' przekazywaną przez utwór.
  Utwory o wysokiej walencji brzmią bardziej radośnie, euforycznie,
  natomiast o niskiej walencji -- smutno, depresyjnie lub gniewnie.
\end{itemize}

\begin{enumerate}
\def\labelenumi{\arabic{enumi}.}
\setcounter{enumi}{2}
\tightlist
\item
  Cechy techniczne i kompozycyjne:
\end{enumerate}

\begin{itemize}
\item
  duration\_ms (numeryczna): Czas trwania utworu wyrażony w
  milisekundach.
\item
  loudness (numeryczna): Ogólna głośność utworu w decybelach (dB),
  uśredniona dla całego nagrania. Zazwyczaj przyjmuje wartości od -60 do
  0 dB.
\item
  tempo (numeryczna): Szacowane tempo utworu w uderzeniach na minutę
  (BPM).
\item
  key (kategoryczna): Tonacja, w jakiej znajduje się utwór, zmapowana
  przy użyciu standardowej notacji Pitch Class (np. 0 = C, 1 = C♯/D♭, 2
  = D).
\item
  mode (kategoryczna/binarna): Tryb skali muzycznej (skala
  majorowa/wesoła = 1, minorowa/smutna = 0).
\item
  time\_signature (numeryczna): Szacowane metrum utworu określające, ile
  uderzeń przypada na jeden takt (najczęściej 4/4).
\end{itemize}

\begin{enumerate}
\def\labelenumi{\arabic{enumi}.}
\setcounter{enumi}{3}
\tightlist
\item
  Metadane i identyfikatory (zmienne pomocnicze / usunięte w
  modelowaniu):
\end{enumerate}

\begin{itemize}
\item
  track\_id: Unikalny identyfikator utworu w systemie Spotify.
\item
  track\_name, artists, album\_name: Tekstowe nazwy utworu, wykonawcy
  oraz albumu.
\item
  track\_genre (kategoryczna): Gatunek muzyczny, do którego przypisany
  jest utwór w ekosystemie Spotify.
\item
  explicit (logiczna): Flaga TRUE / FALSE oznaczająca obecność
  wulgaryzmów lub treści nieodpowiednich dla dzieci.
\end{itemize}

\subsection{Zdefiniowanie problemu predykcyjnego oraz scharakteryzowanie
analizowanego zbioru
danych}\label{zdefiniowanie-problemu-predykcyjnego-oraz-scharakteryzowanie-analizowanego-zbioru-danych}

\textbf{Problem predykcyjny:} Głównym celem analitycznym projektu jest
rozwiązanie problemu klasyfikacji binarnej polegającego na
przewidywaniu, czy dany utwór muzyczny osiągnie status ``hitu''. Na
potrzeby tej analizy zdefiniowano, że utwór staje się hitem (zmienna
docelowa przyjmuje wartość 1), gdy jego wskaźnik \texttt{popularity}
(popularność obliczana przez algorytmy Spotify w skali od 0 do 100) jest
większy lub równy 75. W przeciwnym razie utwór klasyfikowany jest jako
``brak hitu'' (wartość 0).

\textbf{Charakterystyka zbioru:} Zbiór charakteryzuje się występowaniem
silnego niezbalansowania klas (ang. \emph{class imbalance}), ponieważ
utworów o bardzo wysokiej popularności jest w populacji znacznie mniej
niż utworów przeciętnych. Dane zawierają zmienne numeryczne ciągłe (np.
tempo, loudness), zmienne kategoryczne (np. track\_genre, key) oraz
zmienne logiczne (np. explicit). Wymaga to odpowiedniego przygotowania,
w tym kodowania zmiennych jakościowych, skalowania oraz potencjalnie
balansowania klas przed etapem uczenia modeli.

\subsection{Przegląd literatury dotyczącej analizowanego problemu
predykcyjnego oraz omówienie dotychczasowych osiągnięć w tym
zakresie}\label{przeglux105d-literatury-dotyczux105cej-analizowanego-problemu-predykcyjnego-oraz-omuxf3wienie-dotychczasowych-osiux105gniux119ux107-w-tym-zakresie}

Problem przewidywania sukcesu komercyjnego utworów muzycznych jest
często eksplorowany w literaturze z zakresu \emph{Music Information
Retrieval (MIR)} oraz w projektach z dziedziny Data Science.

Dotychczasowe badania wskazują na to, że:

\begin{itemize}
\item
  Cechy psychoakustyczne, takie jak \texttt{danceability} (taneczność),
  \texttt{energy} (energia) oraz \texttt{loudness} (głośność), wykazują
  często zauważalną korelację z popularnością utworów w głównym nurcie
  muzycznym. Wczesne prace nad automatyczną predykcją hitów, takie jak
  badanie Dhanaraj \& Logan {[}1{]}, potwierdzają przydatność cech
  spektralnych i rytmicznych w odróżnianiu piosenek komercyjnie udanych
  od pozostałych.
\item
  Zmienna \texttt{acousticness} (akustyczność) jest zazwyczaj ujemnie
  skorelowana z komercyjnym sukcesem we współczesnej muzyce popularnej.
  Interiano i in. {[}2{]} w analizie obejmującej dekady nagrań wykazali,
  że trendy muzyczne ewoluują w kierunku bardziej ``wyprodukowanego'',
  elektronicznego brzmienia, co potwierdzają dane Spotify.
\item
  Wielu badaczy wskazuje, że proste modele liniowe (jak Regresja
  Logistyczna) często ustępują miejsca modelom opartym na zespołach
  drzew decyzyjnych (np. Random Forest, Gradient Boosting, XGBoost
  {[}3{]}), które znacznie lepiej radzą sobie ze złożonymi, nieliniowymi
  interakcjami między cechami akustycznymi a gustami słuchaczy.
\item
  Osiągnięcie bardzo wysokiej skuteczności (np. F1-score) jest trudne,
  ponieważ na popularność piosenki wpływają również czynniki zewnętrzne,
  których nie ma w zbiorze. Już Pachet \& Roy {[}4{]} w swoim
  przełomowym artykule wykazali, że ``Hit Song Science'' w oparciu
  wyłącznie o cechy audio jest problemem fundamentalnie trudnym ze
  względu na rolę marketingu, artysty i przypadku.
\item
  Problem niezbalansowania klas, będący centralnym wyzwaniem niniejszego
  projektu, jest rozwiązywany w literaturze m.in. za pomocą algorytmu
  SMOTE zaproponowanego przez Chawla i in. {[}5{]}, który generuje
  syntetyczne obserwacje klasy mniejszościowej poprzez interpolację w
  przestrzeni cech.
\end{itemize}

\subsection{Czyszczenie danych oraz wstępna selekcja
cech}\label{czyszczenie-danych-oraz-wstux119pna-selekcja-cech}

Proces przygotowania zbioru do analizy rozpoczął się od czyszczenia
danych i wstępnej selekcji cech. Zrealizowano następujące kroki: *
\textbf{Tworzenie zmiennej objaśnianej:} Na podstawie ustalonego progu
wygenerowano nową, binarną zmienną docelową \texttt{is\_hit}. *
\textbf{Usuwanie duplikatów i braków:} Zbiór Spotify często zawiera te
same utwory przypisane do różnych albumów lub rynków. Zastosowano
filtrowanie pozostawiające tylko unikalne pary tytułu utworu i artysty.
Następnie usunięto wszystkie wiersze zawierające braki danych (NA), aby
zapewnić stabilność modeli. * \textbf{Eliminacja wycieku danych (Data
Leakage):} Kluczowym krokiem było usunięcie oryginalnej zmiennej
\texttt{popularity}. Pozostawienie jej w zbiorze treningowym sprawiłoby,
że model ``oszukiwałby'', mając bezpośredni dostęp do informacji, z
której wyliczono zmienną docelową. * \textbf{Selekcja cech:} Usunięto
kolumny pełniące wyłącznie funkcję identyfikatorów (np. indeksy z pliku
CSV, \texttt{track\_id}), ponieważ nie niosą one żadnej wartości
predykcyjnej dla algorytmów uczenia maszynowego. Po procesie
deduplikacji usunięto również zmienne czysto tekstowe
(\texttt{track\_name}, \texttt{artists}, \texttt{album\_name}), które w
surowej formie nie nadają się do standardowego modelowania
tabelarycznego.

\begin{Shaded}
\begin{Highlighting}[]
\FunctionTok{library}\NormalTok{(dplyr)}
\end{Highlighting}
\end{Shaded}

\begin{verbatim}
## Warning: pakiet 'dplyr' został zbudowany w wersji R 4.4.2
\end{verbatim}

\begin{verbatim}
## 
## Dołączanie pakietu: 'dplyr'
\end{verbatim}

\begin{verbatim}
## Następujące obiekty zostały zakryte z 'package:stats':
## 
##     filter, lag
\end{verbatim}

\begin{verbatim}
## Następujące obiekty zostały zakryte z 'package:base':
## 
##     intersect, setdiff, setequal, union
\end{verbatim}

\begin{Shaded}
\begin{Highlighting}[]
\FunctionTok{library}\NormalTok{(readr)}
\FunctionTok{library}\NormalTok{(tidyr)}

\NormalTok{df }\OtherTok{\textless{}{-}} \FunctionTok{read\_csv}\NormalTok{(}\StringTok{"spotify{-}tracks{-}dataset.csv"}\NormalTok{)}
\end{Highlighting}
\end{Shaded}

\begin{verbatim}
## New names:
## * `` -> `...1`
\end{verbatim}

\begin{verbatim}
## Rows: 114000 Columns: 22
## -- Column specification --------------------------------------------------------
## Delimiter: ","
## chr  (5): track_id, artists, album_name, track_name, track_genre
## dbl (16): ...1, Unnamed: 0, popularity, duration_ms, danceability, energy, k...
## lgl  (1): explicit
## 
## i Use `spec()` to retrieve the full column specification for this data.
## i Specify the column types or set `show_col_types = FALSE` to quiet this message.
\end{verbatim}

\begin{Shaded}
\begin{Highlighting}[]
\NormalTok{POPULARITY\_THRESHOLD }\OtherTok{\textless{}{-}} \DecValTok{75}

\NormalTok{df }\OtherTok{\textless{}{-}}\NormalTok{ df }\SpecialCharTok{\%\textgreater{}\%}
  \FunctionTok{mutate}\NormalTok{(}\AttributeTok{is\_hit =} \FunctionTok{as.factor}\NormalTok{(}\FunctionTok{ifelse}\NormalTok{(popularity }\SpecialCharTok{\textgreater{}=}\NormalTok{ POPULARITY\_THRESHOLD, }\DecValTok{1}\NormalTok{, }\DecValTok{0}\NormalTok{)))}

\NormalTok{df }\OtherTok{\textless{}{-}}\NormalTok{ df }\SpecialCharTok{\%\textgreater{}\%}
  \FunctionTok{distinct}\NormalTok{(track\_name, artists, }\AttributeTok{.keep\_all =} \ConstantTok{TRUE}\NormalTok{) }\SpecialCharTok{\%\textgreater{}\%}
  \FunctionTok{drop\_na}\NormalTok{()}

\NormalTok{df }\OtherTok{\textless{}{-}}\NormalTok{ df }\SpecialCharTok{\%\textgreater{}\%}
  \FunctionTok{select}\NormalTok{(}\SpecialCharTok{{-}}\FunctionTok{any\_of}\NormalTok{(}\FunctionTok{c}\NormalTok{(}\StringTok{"...1"}\NormalTok{, }\StringTok{"Unnamed: 0"}\NormalTok{, }\StringTok{"X1"}\NormalTok{, }\StringTok{"track\_id"}\NormalTok{, }\StringTok{"popularity"}\NormalTok{, }\StringTok{"track\_name"}\NormalTok{, }\StringTok{"artists"}\NormalTok{, }\StringTok{"album\_name"}\NormalTok{)))}
\end{Highlighting}
\end{Shaded}

\textbf{Wniosek:} Po przeprowadzeniu czyszczenia uzyskano spójny,
pozbawiony duplikatów i braków zbiór danych, gotowy do dalszej analizy.
Eliminacja zmiennej \texttt{popularity} jest krytycznym krokiem
zapobiegającym wyciekowi danych (data leakage) -- bez tego modelu
trenowałby się de facto na odpowiedzi. Usunięcie kolumn tekstowych i
identyfikatorów pozwala skupić się wyłącznie na cechach o potencjale
predykcyjnym. Wynikowy zbiór zawiera zarówno zmienne numeryczne
(parametry audio), jak i kategoryczne (\texttt{track\_genre},
\texttt{key}, \texttt{mode}), co determinuje dalsze kroki przetwarzania
w pipeline'u modelowania.

\subsection{Omówienie oraz zwizualizowanie istotnych zależności w
analizowanym zbiorze
danych}\label{omuxf3wienie-oraz-zwizualizowanie-istotnych-zaleux17cnoux15bci-w-analizowanym-zbiorze-danych}

W celu lepszego zrozumienia struktury danych oraz relacji między
zmiennymi przeprowadzono eksploracyjną analizę danych (EDA) z
wykorzystaniem technik wizualizacji. Skupiono się na trzech kluczowych
aspektach:

\textbf{1. Rozkład zmiennej docelowej:} Wykres słupkowy ilustrujący
rozkład zmiennej \texttt{is\_hit} ukazuje drastyczne niezbalansowanie
klas. Zdecydowana większość obserwacji w zbiorze należy do klasy
większościowej (0 - ``brak hitu''), podczas gdy utwory zaklasyfikowane
jako hity (1) stanowią jedynie niewielki ułamek populacji (ok. 5-7\%).

\emph{Wniosek:} Skrajna asymetria klas sugeruje, że model nauczony na
surowych danych mógłby osiągnąć wysoką dokładność (Accuracy) poprzez
proste przypisywanie wszystkim utworom klasy ``0''. Wymusza to w
późniejszych etapach zastosowanie technik balansowania (SMOTE) oraz
metryk oceny skupionych na klasie mniejszościowej (Precision, Recall,
F1-score).

\textbf{2. Macierz korelacji:} Analiza korelogramu dla zmiennych
numerycznych ujawnia kilka silnych zależności, które są spójne z
intuicją muzyczną. Występuje bardzo silna, dodatnia korelacja pomiędzy
zmiennymi \texttt{energy} (energia) oraz \texttt{loudness} (głośność).
Utwory głośniejsze są przez algorytmy Spotify klasyfikowane jako
bardziej energetyczne. Zmienna \texttt{acousticness} (akustyczność) jest
silnie ujemnie skorelowana zarówno z energią, jak i głośnością, co
odzwierciedla techniczną różnicę między instrumentami akustycznymi a
zelektronizowaną produkcją współczesną.

\emph{Wniosek:} Identyfikacja silnie skorelowanych par zmiennych
objaśniających (multikolinearność) jest sygnałem, że proste modele
liniowe mogą mieć trudności z poprawną estymacją wag. W procesie
modelowania preferowane będą algorytmy odporne na to zjawisko, takie jak
lasy losowe, lub konieczne będzie usunięcie redundantnych cech.

\textbf{3. Redukcja wymiarów (PCA):} Aby zwizualizować trudność
postawionego problemu analitycznego, przeprowadzono Analizę Głównych
Składowych (PCA) dla numerycznych cech akustycznych. Technika ta
pozwoliła na redukcję wielowymiarowej przestrzeni danych do dwóch
głównych wymiarów (PC1 i PC2), które łącznie wyjaśniają znaczną część
wariancji zbioru.

\emph{Wniosek:} Nałożenie na wykres punktowy kolorów odpowiadających
klasie docelowej uwidacznia, że klasa ``hitów'' nie tworzy na rynku
muzycznym jednej, wyraźnie odseparowanej grupy, lecz jest głęboko
zanurzona w przestrzeni cech utworów przeciętnych. Analiza ładunków
sugeruje, że PC1 jest silnie związana z ``intensywnością'' utworu
(energia, głośność vs akustyczność), natomiast PC2 może odzwierciedlać
parametry rytmiczne i nastrojowe (taneczność, walencja). Fakt, że hity
są rozproszone, oznacza, że nie istnieje prosty zestaw reguł (``złota
formuła''), który gwarantuje sukces. Granica decyzyjna jest wysoce
nieliniowa, co uzasadnia użycie modeli takich jak XGBoost czy Random
Forest.

\begin{Shaded}
\begin{Highlighting}[]
\FunctionTok{library}\NormalTok{(ggplot2)}
\end{Highlighting}
\end{Shaded}

\begin{verbatim}
## Warning: pakiet 'ggplot2' został zbudowany w wersji R 4.4.3
\end{verbatim}

\begin{Shaded}
\begin{Highlighting}[]
\FunctionTok{library}\NormalTok{(corrplot)}
\end{Highlighting}
\end{Shaded}

\begin{verbatim}
## Warning: pakiet 'corrplot' został zbudowany w wersji R 4.4.3
\end{verbatim}

\begin{verbatim}
## corrplot 0.95 loaded
\end{verbatim}

\begin{Shaded}
\begin{Highlighting}[]
\FunctionTok{library}\NormalTok{(tidymodels)}
\end{Highlighting}
\end{Shaded}

\begin{verbatim}
## Warning: pakiet 'tidymodels' został zbudowany w wersji R 4.4.3
\end{verbatim}

\begin{verbatim}
## -- Attaching packages -------------------------------------- tidymodels 1.4.1 --
\end{verbatim}

\begin{verbatim}
## v broom        1.0.12     v rsample      1.3.2 
## v dials        1.4.2      v tailor       0.1.0 
## v infer        1.1.0      v tune         2.0.1 
## v modeldata    1.5.1      v workflows    1.3.0 
## v parsnip      1.4.1      v workflowsets 1.1.1 
## v purrr        1.2.1      v yardstick    1.3.2 
## v recipes      1.3.1
\end{verbatim}

\begin{verbatim}
## Warning: pakiet 'broom' został zbudowany w wersji R 4.4.3
\end{verbatim}

\begin{verbatim}
## Warning: pakiet 'dials' został zbudowany w wersji R 4.4.3
\end{verbatim}

\begin{verbatim}
## Warning: pakiet 'scales' został zbudowany w wersji R 4.4.3
\end{verbatim}

\begin{verbatim}
## Warning: pakiet 'infer' został zbudowany w wersji R 4.4.3
\end{verbatim}

\begin{verbatim}
## Warning: pakiet 'modeldata' został zbudowany w wersji R 4.4.3
\end{verbatim}

\begin{verbatim}
## Warning: pakiet 'parsnip' został zbudowany w wersji R 4.4.3
\end{verbatim}

\begin{verbatim}
## Warning: pakiet 'purrr' został zbudowany w wersji R 4.4.3
\end{verbatim}

\begin{verbatim}
## Warning: pakiet 'recipes' został zbudowany w wersji R 4.4.3
\end{verbatim}

\begin{verbatim}
## Warning: pakiet 'rsample' został zbudowany w wersji R 4.4.3
\end{verbatim}

\begin{verbatim}
## Warning: pakiet 'tailor' został zbudowany w wersji R 4.4.3
\end{verbatim}

\begin{verbatim}
## Warning: pakiet 'tune' został zbudowany w wersji R 4.4.3
\end{verbatim}

\begin{verbatim}
## Warning: pakiet 'workflows' został zbudowany w wersji R 4.4.3
\end{verbatim}

\begin{verbatim}
## Warning: pakiet 'workflowsets' został zbudowany w wersji R 4.4.3
\end{verbatim}

\begin{verbatim}
## Warning: pakiet 'yardstick' został zbudowany w wersji R 4.4.3
\end{verbatim}

\begin{verbatim}
## -- Conflicts ----------------------------------------- tidymodels_conflicts() --
## x purrr::discard()  masks scales::discard()
## x dplyr::filter()   masks stats::filter()
## x dplyr::lag()      masks stats::lag()
## x yardstick::spec() masks readr::spec()
## x recipes::step()   masks stats::step()
\end{verbatim}

\begin{Shaded}
\begin{Highlighting}[]
\FunctionTok{ggplot}\NormalTok{(df, }\FunctionTok{aes}\NormalTok{(}\AttributeTok{x =}\NormalTok{ is\_hit, }\AttributeTok{fill =}\NormalTok{ is\_hit)) }\SpecialCharTok{+}
  \FunctionTok{geom\_bar}\NormalTok{(}\AttributeTok{show.legend =} \ConstantTok{FALSE}\NormalTok{) }\SpecialCharTok{+}
  \FunctionTok{scale\_fill\_manual}\NormalTok{(}\AttributeTok{values =} \FunctionTok{c}\NormalTok{(}\StringTok{"0"} \OtherTok{=} \StringTok{"steelblue"}\NormalTok{, }\StringTok{"1"} \OtherTok{=} \StringTok{"firebrick"}\NormalTok{)) }\SpecialCharTok{+}
  \FunctionTok{theme\_minimal}\NormalTok{() }\SpecialCharTok{+}
  \FunctionTok{labs}\NormalTok{(}\AttributeTok{title =} \StringTok{"Rozkład zmiennej docelowej (is\_hit)"}\NormalTok{,}
       \AttributeTok{x =} \StringTok{"Klasa (0 = Brak hitu, 1 = Hit)"}\NormalTok{,}
       \AttributeTok{y =} \StringTok{"Liczba utworów"}\NormalTok{)}
\end{Highlighting}
\end{Shaded}

\pandocbounded{\includegraphics[keepaspectratio]{projekt-final_files/figure-latex/unnamed-chunk-2-1.pdf}}

\textbf{Wniosek:} Wykres potwierdza silną asymetrię klas -- hity (klasa
1) stanowią zdecydowaną mniejszość w zbiorze, co sprawia, że model
uczony na surowych danych mógłby niemal bez wysiłku uzyskać wysoką
dokładność ogólną, jednocześnie całkowicie ignorując klasę
mniejszościową. Uzasadnia to konieczność zastosowania techniki SMOTE w
dalszym etapie modelowania oraz skupienia się na metrykach Precision,
Recall i F1-Score zamiast samej dokładności.

\begin{Shaded}
\begin{Highlighting}[]
\FunctionTok{library}\NormalTok{(dplyr)}
\FunctionTok{library}\NormalTok{(corrplot)}

\NormalTok{numeric\_cols }\OtherTok{\textless{}{-}}\NormalTok{ df }\SpecialCharTok{\%\textgreater{}\%}
  \FunctionTok{select}\NormalTok{(}\FunctionTok{where}\NormalTok{(is.numeric))}

\NormalTok{corr\_matrix }\OtherTok{\textless{}{-}} \FunctionTok{cor}\NormalTok{(numeric\_cols, }\AttributeTok{use =} \StringTok{"complete.obs"}\NormalTok{)}

\FunctionTok{corrplot}\NormalTok{(corr\_matrix, }
         \AttributeTok{method =} \StringTok{"color"}\NormalTok{, }
         \AttributeTok{type =} \StringTok{"upper"}\NormalTok{, }
         \AttributeTok{tl.col =} \StringTok{"black"}\NormalTok{, }
         \AttributeTok{tl.cex =} \FloatTok{0.8}\NormalTok{)}
\end{Highlighting}
\end{Shaded}

\pandocbounded{\includegraphics[keepaspectratio]{projekt-final_files/figure-latex/unnamed-chunk-3-1.pdf}}

\textbf{Wniosek:} Macierz korelacji ujawnia kilka kluczowych zależności
strukturalnych. Najsilniejsza jest dodatnia korelacja między
\texttt{energy} a \texttt{loudness} (r ≈ 0.76), co potwierdza, że
głośność jest technicznym komponentem postrzeganej energii. Silna ujemna
korelacja \texttt{acousticness} z oboma tymi zmiennymi wskazuje, że
instrumenty akustyczne i elektroniczna produkcja muzyczna zajmują
przeciwne bieguny przestrzeni cech. Z perspektywy modelowania, wysoka
współliniowość może destabilizować estymację wag w modelu regresji
logistycznej -- kolejny argument przemawiający za algorytmami
drzewiastymi, odpornymi na to zjawisko.

\begin{Shaded}
\begin{Highlighting}[]
\FunctionTok{library}\NormalTok{(ggplot2)}
\FunctionTok{library}\NormalTok{(dplyr)}
\FunctionTok{library}\NormalTok{(tidymodels)}

\NormalTok{pca\_rec }\OtherTok{\textless{}{-}} \FunctionTok{recipe}\NormalTok{(}\SpecialCharTok{\textasciitilde{}}\NormalTok{ ., }\AttributeTok{data =}\NormalTok{ df }\SpecialCharTok{\%\textgreater{}\%} \FunctionTok{select}\NormalTok{(}\FunctionTok{where}\NormalTok{(is.numeric), is\_hit)) }\SpecialCharTok{\%\textgreater{}\%}
  \FunctionTok{step\_normalize}\NormalTok{(}\FunctionTok{all\_numeric\_predictors}\NormalTok{()) }\SpecialCharTok{\%\textgreater{}\%}
  \FunctionTok{step\_pca}\NormalTok{(}\FunctionTok{all\_numeric\_predictors}\NormalTok{(), }\AttributeTok{num\_comp =} \DecValTok{2}\NormalTok{)}

\NormalTok{pca\_prep }\OtherTok{\textless{}{-}} \FunctionTok{prep}\NormalTok{(pca\_rec)}
\NormalTok{pca\_data }\OtherTok{\textless{}{-}} \FunctionTok{bake}\NormalTok{(pca\_prep, }\AttributeTok{new\_data =} \ConstantTok{NULL}\NormalTok{)}

\FunctionTok{ggplot}\NormalTok{(pca\_data, }\FunctionTok{aes}\NormalTok{(}\AttributeTok{x =}\NormalTok{ PC1, }\AttributeTok{y =}\NormalTok{ PC2, }\AttributeTok{color =}\NormalTok{ is\_hit)) }\SpecialCharTok{+}
  \FunctionTok{geom\_point}\NormalTok{(}\AttributeTok{alpha =} \FloatTok{0.6}\NormalTok{, }\AttributeTok{size =} \FloatTok{1.5}\NormalTok{) }\SpecialCharTok{+}
  \FunctionTok{scale\_color\_manual}\NormalTok{(}\AttributeTok{values =} \FunctionTok{c}\NormalTok{(}\StringTok{"0"} \OtherTok{=} \StringTok{"steelblue"}\NormalTok{, }\StringTok{"1"} \OtherTok{=} \StringTok{"firebrick"}\NormalTok{)) }\SpecialCharTok{+}
  \FunctionTok{theme\_minimal}\NormalTok{() }\SpecialCharTok{+}
  \FunctionTok{labs}\NormalTok{(}\AttributeTok{title =} \StringTok{"Redukcja wymiarów (PCA) {-} Dystrybucja hitów na rynku"}\NormalTok{,}
       \AttributeTok{x =} \StringTok{"Główna Składowa 1 (PC1)"}\NormalTok{,}
       \AttributeTok{y =} \StringTok{"Główna Składowa 2 (PC2)"}\NormalTok{,}
       \AttributeTok{color =} \StringTok{"Hit"}\NormalTok{)}
\end{Highlighting}
\end{Shaded}

\pandocbounded{\includegraphics[keepaspectratio]{projekt-final_files/figure-latex/unnamed-chunk-4-1.pdf}}

\textbf{Wniosek:} Wykres PCA obrazuje fundamentalną trudność
postawionego problemu -- klasa hitów (czerwone punkty) nie tworzy
zwartego, separowalnego skupiska, lecz jest rozproszona w całej
przestrzeni cech. Oznacza to, że nie istnieje prosta, linearna granica
decyzyjna mogąca skutecznie oddzielić hity od pozostałych utworów. Układ
PC1 (intensywność produkcji) i PC2 (parametry rytmiczno-nastrojowe)
wspólnie wyjaśniają główną zmienność zbioru, jednak nie wystarczają do
jednoznacznej klasyfikacji. Ten wynik jednoznacznie uzasadnia
zastosowanie modeli o nieliniowej granicy decyzyjnej w dalszych etapach
analizy.

\subsection{Przygotowanie danych do dalszego procesu
modelowania}\label{przygotowanie-danych-do-dalszego-procesu-modelowania}

Przygotowanie danych do uczenia maszynowego zrealizowano z
wykorzystaniem ekosystemu tidymodels. Proces ten podzielono na
następujące fazy:

Podział zbioru: Dane podzielono na zbiór treningowy (80\% obserwacji)
oraz testowy (20\% obserwacji). Zastosowano podział warstwowy
(stratified split) ze względu na zmienną is\_hit, co gwarantuje, że
proporcja hitów w obu zbiorach pozostanie identyczna jak w populacji
bazowej.

Transformacja predyktorów: Utworzono recepturę transformacji (recipe), w
której zastosowano kodowanie ``gorącojedynkowe'' (One-Hot Encoding /
Dummy Variables) dla wszystkich predyktorów kategorycznych (np.
track\_genre) oraz standaryzację (centrowanie i skalowanie) dla
wszystkich zmiennych numerycznych. Zmienna logiczna explicit została
wcześniej przekształcona na typ numeryczny, aby zapobiec błędom
przetwarzania geometrycznego.

Balansowanie klas: Ze względu na zidentyfikowane we wcześniejszym etapie
skrajne niezbalansowanie klas, do receptury włączono krok oversamplingu
przy użyciu algorytmu SMOTE (Synthetic Minority Over-sampling
Technique). Został on zastosowany wyłącznie na zbiorze treningowym, aby
wygenerować syntetyczne obserwacje klasy mniejszościowej (hitów) i
zrównoważyć zbiór, co zapobiegnie faworyzowaniu klasy większościowej
przez modele.

\begin{Shaded}
\begin{Highlighting}[]
\FunctionTok{library}\NormalTok{(tidymodels)}
\FunctionTok{library}\NormalTok{(themis)}
\end{Highlighting}
\end{Shaded}

\begin{verbatim}
## Warning: pakiet 'themis' został zbudowany w wersji R 4.4.3
\end{verbatim}

\begin{Shaded}
\begin{Highlighting}[]
\NormalTok{df }\OtherTok{\textless{}{-}}\NormalTok{ df }\SpecialCharTok{\%\textgreater{}\%}
  \FunctionTok{mutate}\NormalTok{(}\AttributeTok{explicit =} \FunctionTok{as.numeric}\NormalTok{(explicit))}

\FunctionTok{set.seed}\NormalTok{(}\DecValTok{42}\NormalTok{)}

\NormalTok{data\_split }\OtherTok{\textless{}{-}} \FunctionTok{initial\_split}\NormalTok{(df, }\AttributeTok{prop =} \FloatTok{0.8}\NormalTok{, }\AttributeTok{strata =}\NormalTok{ is\_hit)}
\NormalTok{train\_data }\OtherTok{\textless{}{-}} \FunctionTok{training}\NormalTok{(data\_split)}
\NormalTok{test\_data }\OtherTok{\textless{}{-}} \FunctionTok{testing}\NormalTok{(data\_split)}

\NormalTok{rec }\OtherTok{\textless{}{-}} \FunctionTok{recipe}\NormalTok{(is\_hit }\SpecialCharTok{\textasciitilde{}}\NormalTok{ ., }\AttributeTok{data =}\NormalTok{ train\_data) }\SpecialCharTok{\%\textgreater{}\%}
  \FunctionTok{step\_dummy}\NormalTok{(}\FunctionTok{all\_nominal\_predictors}\NormalTok{()) }\SpecialCharTok{\%\textgreater{}\%}
  \FunctionTok{step\_normalize}\NormalTok{(}\FunctionTok{all\_numeric\_predictors}\NormalTok{()) }\SpecialCharTok{\%\textgreater{}\%}
  \FunctionTok{step\_smote}\NormalTok{(is\_hit)}

\NormalTok{prep\_rec }\OtherTok{\textless{}{-}} \FunctionTok{prep}\NormalTok{(rec)}

\NormalTok{train\_processed }\OtherTok{\textless{}{-}} \FunctionTok{bake}\NormalTok{(prep\_rec, }\AttributeTok{new\_data =} \ConstantTok{NULL}\NormalTok{)}
\NormalTok{test\_processed }\OtherTok{\textless{}{-}} \FunctionTok{bake}\NormalTok{(prep\_rec, }\AttributeTok{new\_data =}\NormalTok{ test\_data)}

\FunctionTok{dim}\NormalTok{(train\_processed)}
\end{Highlighting}
\end{Shaded}

\begin{verbatim}
## [1] 128074    127
\end{verbatim}

\begin{Shaded}
\begin{Highlighting}[]
\FunctionTok{table}\NormalTok{(train\_processed}\SpecialCharTok{$}\NormalTok{is\_hit)}
\end{Highlighting}
\end{Shaded}

\begin{verbatim}
## 
##     0     1 
## 64037 64037
\end{verbatim}

\begin{Shaded}
\begin{Highlighting}[]
\FunctionTok{head}\NormalTok{(train\_processed)}
\end{Highlighting}
\end{Shaded}

\begin{verbatim}
## # A tibble: 6 x 127
##   duration_ms explicit danceability energy    key loudness   mode speechiness
##         <dbl>    <dbl>        <dbl>  <dbl>  <dbl>    <dbl>  <dbl>       <dbl>
## 1    -0.00400   -0.306        0.657 -0.673 -1.20    0.349  -1.32        0.463
## 2    -0.727     -0.306       -0.784 -1.81  -1.20   -1.63    0.759      -0.108
## 3    -0.260     -0.306       -1.65  -2.23  -1.48   -1.88    0.759      -0.450
## 4    -0.151     -0.306        0.724 -0.595  0.203  -0.0405  0.759       0.138
## 5    -0.0153    -0.306       -0.857 -1.89  -0.921  -0.0433  0.759      -0.457
## 6     0.106     -0.306        0.809 -0.739  1.61   -0.139   0.759      -0.404
## # i 119 more variables: acousticness <dbl>, instrumentalness <dbl>,
## #   liveness <dbl>, valence <dbl>, tempo <dbl>, time_signature <dbl>,
## #   track_genre_afrobeat <dbl>, track_genre_alt.rock <dbl>,
## #   track_genre_alternative <dbl>, track_genre_ambient <dbl>,
## #   track_genre_anime <dbl>, track_genre_black.metal <dbl>,
## #   track_genre_bluegrass <dbl>, track_genre_blues <dbl>,
## #   track_genre_brazil <dbl>, track_genre_breakbeat <dbl>, ...
\end{verbatim}

\textbf{Wniosek:} Po zastosowaniu SMOTE zbiór treningowy zawiera
wyrównany rozkład klas -- liczba syntetycznych obserwacji klasy ``hit''
wzrosła do poziomu klasy większościowej. Podział warstwowy (stratified
split) zapewnia, że proporcja hitów w zbiorze testowym odpowiada
proporcji w oryginalnym zbiorze, co gwarantuje wiarygodną ewaluację.
Standaryzacja zmiennych numerycznych jest niezbędna dla stabilności
regresji logistycznej i przyspiesza zbieżność XGBoost. Kodowanie One-Hot
Encoding zmiennej \texttt{track\_genre} znacznie zwiększa liczbę
predyktorów, ale pozwala modelom drzewiastym efektywnie wykorzystać
informację o gatunku muzycznym. Zbiór testowy pozostaje celowo
niezmieniony (bez SMOTE), aby ocena odzwierciedlała realne,
niezbalansowane warunki.

\subsection{Budowa, trenowanie oraz wstępne porównanie wybranych modeli
uczenia
maszynowego}\label{budowa-trenowanie-oraz-wstux119pne-poruxf3wnanie-wybranych-modeli-uczenia-maszynowego}

W procesie modelowania predykcyjnego zdecydowano się na wykorzystanie i
porównanie trzech algorytmów klasyfikacyjnych o zróżnicowanej
charakterystyce matematycznej, co pozwala na rzetelną ocenę
skuteczności:

Regresja Logistyczna (Logistic Regression): Zastosowana jako model
bazowy (tzw. baseline). Jest to model parametryczny, łatwy w
interpretacji, jednak zakładający liniową zależność między predyktorami
a logarytmem szansy wystąpienia klasy pozytywnej.

Las Losowy (Random Forest): Zaawansowany model zespołowy (ang.
ensemble), oparty na agregacji wyników wielu drzew decyzyjnych w ramach
techniki Baggingu. Wybór tego modelu podyktowany był jego wysoką
skutecznością, brakiem wrażliwości na zjawisko multikolinearności oraz
zdolnością do modelowania złożonych, nieliniowych relacji.

XGBoost (Extreme Gradient Boosting): Algorytm oparty na technice
Boostingu (budowa drzew sekwencyjna, optymalizująca błędy poprzedników),
będący obecnie jednym z najskuteczniejszych algorytmów klasyfikacyjnych
dla strukturalnych danych tabelarycznych.

Modele zostały wytrenowane na przetworzonym i zbalansowanym zbiorze
treningowym. Następnie wygenerowano predykcje na niewidzianym wcześniej
zbiorze testowym, co stanowi fundament do obiektywnej ewaluacji w
kolejnym kroku.

\begin{Shaded}
\begin{Highlighting}[]
\FunctionTok{library}\NormalTok{(parsnip)}
\FunctionTok{library}\NormalTok{(ranger)}
\end{Highlighting}
\end{Shaded}

\begin{verbatim}
## Warning: pakiet 'ranger' został zbudowany w wersji R 4.4.2
\end{verbatim}

\begin{Shaded}
\begin{Highlighting}[]
\FunctionTok{library}\NormalTok{(xgboost)}
\end{Highlighting}
\end{Shaded}

\begin{verbatim}
## Warning: pakiet 'xgboost' został zbudowany w wersji R 4.4.3
\end{verbatim}

\begin{Shaded}
\begin{Highlighting}[]
\NormalTok{log\_spec }\OtherTok{\textless{}{-}} \FunctionTok{logistic\_reg}\NormalTok{() }\SpecialCharTok{\%\textgreater{}\%}
  \FunctionTok{set\_engine}\NormalTok{(}\StringTok{"glm"}\NormalTok{) }\SpecialCharTok{\%\textgreater{}\%}
  \FunctionTok{set\_mode}\NormalTok{(}\StringTok{"classification"}\NormalTok{)}

\NormalTok{rf\_spec }\OtherTok{\textless{}{-}} \FunctionTok{rand\_forest}\NormalTok{(}\AttributeTok{trees =} \DecValTok{100}\NormalTok{) }\SpecialCharTok{\%\textgreater{}\%}
  \FunctionTok{set\_engine}\NormalTok{(}\StringTok{"ranger"}\NormalTok{, }\AttributeTok{importance =} \StringTok{"impurity"}\NormalTok{) }\SpecialCharTok{\%\textgreater{}\%}
  \FunctionTok{set\_mode}\NormalTok{(}\StringTok{"classification"}\NormalTok{)}

\NormalTok{xgb\_spec }\OtherTok{\textless{}{-}} \FunctionTok{boost\_tree}\NormalTok{(}\AttributeTok{trees =} \DecValTok{100}\NormalTok{) }\SpecialCharTok{\%\textgreater{}\%}
  \FunctionTok{set\_engine}\NormalTok{(}\StringTok{"xgboost"}\NormalTok{) }\SpecialCharTok{\%\textgreater{}\%}
  \FunctionTok{set\_mode}\NormalTok{(}\StringTok{"classification"}\NormalTok{)}

\NormalTok{log\_fit }\OtherTok{\textless{}{-}}\NormalTok{ log\_spec }\SpecialCharTok{\%\textgreater{}\%}
  \FunctionTok{fit}\NormalTok{(is\_hit }\SpecialCharTok{\textasciitilde{}}\NormalTok{ ., }\AttributeTok{data =}\NormalTok{ train\_processed)}

\NormalTok{rf\_fit }\OtherTok{\textless{}{-}}\NormalTok{ rf\_spec }\SpecialCharTok{\%\textgreater{}\%}
  \FunctionTok{fit}\NormalTok{(is\_hit }\SpecialCharTok{\textasciitilde{}}\NormalTok{ ., }\AttributeTok{data =}\NormalTok{ train\_processed)}

\NormalTok{xgb\_fit }\OtherTok{\textless{}{-}}\NormalTok{ xgb\_spec }\SpecialCharTok{\%\textgreater{}\%}
  \FunctionTok{fit}\NormalTok{(is\_hit }\SpecialCharTok{\textasciitilde{}}\NormalTok{ ., }\AttributeTok{data =}\NormalTok{ train\_processed)}

\NormalTok{log\_preds }\OtherTok{\textless{}{-}} \FunctionTok{predict}\NormalTok{(log\_fit, }\AttributeTok{new\_data =}\NormalTok{ test\_processed) }\SpecialCharTok{\%\textgreater{}\%}
  \FunctionTok{bind\_cols}\NormalTok{(}\FunctionTok{predict}\NormalTok{(log\_fit, }\AttributeTok{new\_data =}\NormalTok{ test\_processed, }\AttributeTok{type =} \StringTok{"prob"}\NormalTok{)) }\SpecialCharTok{\%\textgreater{}\%}
  \FunctionTok{bind\_cols}\NormalTok{(test\_processed }\SpecialCharTok{\%\textgreater{}\%} \FunctionTok{select}\NormalTok{(is\_hit))}

\NormalTok{rf\_preds }\OtherTok{\textless{}{-}} \FunctionTok{predict}\NormalTok{(rf\_fit, }\AttributeTok{new\_data =}\NormalTok{ test\_processed) }\SpecialCharTok{\%\textgreater{}\%}
  \FunctionTok{bind\_cols}\NormalTok{(}\FunctionTok{predict}\NormalTok{(rf\_fit, }\AttributeTok{new\_data =}\NormalTok{ test\_processed, }\AttributeTok{type =} \StringTok{"prob"}\NormalTok{)) }\SpecialCharTok{\%\textgreater{}\%}
  \FunctionTok{bind\_cols}\NormalTok{(test\_processed }\SpecialCharTok{\%\textgreater{}\%} \FunctionTok{select}\NormalTok{(is\_hit))}

\NormalTok{xgb\_preds }\OtherTok{\textless{}{-}} \FunctionTok{predict}\NormalTok{(xgb\_fit, }\AttributeTok{new\_data =}\NormalTok{ test\_processed) }\SpecialCharTok{\%\textgreater{}\%}
  \FunctionTok{bind\_cols}\NormalTok{(}\FunctionTok{predict}\NormalTok{(xgb\_fit, }\AttributeTok{new\_data =}\NormalTok{ test\_processed, }\AttributeTok{type =} \StringTok{"prob"}\NormalTok{)) }\SpecialCharTok{\%\textgreater{}\%}
  \FunctionTok{bind\_cols}\NormalTok{(test\_processed }\SpecialCharTok{\%\textgreater{}\%} \FunctionTok{select}\NormalTok{(is\_hit))}
\end{Highlighting}
\end{Shaded}

\textbf{Wniosek:} Wszystkie trzy modele zostały pomyślnie wytrenowane na
zbalansowanym zbiorze treningowym. Każdy z nich reprezentuje inną
filozofię uczenia: regresja logistyczna szuka globalnej, liniowej
granicy decyzyjnej; las losowy buduje ensemble niezależnych drzew
(bagging), a XGBoost sekwencyjnie koryguje błędy poprzednich drzew
(boosting). Zróżnicowanie podejść jest celowe -- pozwoli ocenić, które
założenia najlepiej pasują do struktury danych muzycznych. Predykcje na
zbiorze testowym umożliwią obiektywne porównanie modeli w kolejnym
etapie.

\subsubsection{Wstępne porównanie modeli za pomocą kroswalidacji (5-fold
CV)}\label{wstux119pne-poruxf3wnanie-modeli-za-pomocux105-kroswalidacji-5-fold-cv}

Aby uzyskać stabilną i odporną na losowość jednego podziału ocenę
modeli, przeprowadzono \textbf{5-krotną kroswalidację} (\emph{5-fold
Cross-Validation}) na oryginalnym zbiorze treningowym. W każdej z 5
iteracji receptura transformacji (w tym krok SMOTE) jest stosowana
\emph{niezależnie} wyłącznie na danych treningowych danego foldu -- dane
walidacyjne pozostają niezmienione. Takie podejście jest metodologicznie
poprawniejsze niż ocena na jednym podziale, ponieważ uśrednia wyniki po
różnych próbkach i dostarcza estymacji niepewności (odchylenie
standardowe metryki).

\begin{Shaded}
\begin{Highlighting}[]
\FunctionTok{library}\NormalTok{(tidymodels)}
\FunctionTok{library}\NormalTok{(themis)}

\CommentTok{\# Ustawienie klasy "1" jako zdarzenia pozytywnego}
\FunctionTok{options}\NormalTok{(}\AttributeTok{yardstick.event\_first =} \ConstantTok{FALSE}\NormalTok{)}

\FunctionTok{set.seed}\NormalTok{(}\DecValTok{7851}\NormalTok{)}
\NormalTok{cv\_folds }\OtherTok{\textless{}{-}} \FunctionTok{vfold\_cv}\NormalTok{(train\_data, }\AttributeTok{v =} \DecValTok{5}\NormalTok{, }\AttributeTok{strata =}\NormalTok{ is\_hit)}

\CommentTok{\# Tworzenie workflow: recipe (z SMOTE) + model}
\NormalTok{log\_wf }\OtherTok{\textless{}{-}} \FunctionTok{workflow}\NormalTok{() }\SpecialCharTok{\%\textgreater{}\%} \FunctionTok{add\_recipe}\NormalTok{(rec) }\SpecialCharTok{\%\textgreater{}\%} \FunctionTok{add\_model}\NormalTok{(log\_spec)}
\NormalTok{rf\_wf  }\OtherTok{\textless{}{-}} \FunctionTok{workflow}\NormalTok{() }\SpecialCharTok{\%\textgreater{}\%} \FunctionTok{add\_recipe}\NormalTok{(rec) }\SpecialCharTok{\%\textgreater{}\%} \FunctionTok{add\_model}\NormalTok{(rf\_spec)}
\NormalTok{xgb\_wf }\OtherTok{\textless{}{-}} \FunctionTok{workflow}\NormalTok{() }\SpecialCharTok{\%\textgreater{}\%} \FunctionTok{add\_recipe}\NormalTok{(rec) }\SpecialCharTok{\%\textgreater{}\%} \FunctionTok{add\_model}\NormalTok{(xgb\_spec)}

\NormalTok{cv\_met }\OtherTok{\textless{}{-}} \FunctionTok{metric\_set}\NormalTok{(accuracy, precision, recall, f\_meas, roc\_auc)}

\NormalTok{log\_cv }\OtherTok{\textless{}{-}} \FunctionTok{fit\_resamples}\NormalTok{(log\_wf, cv\_folds, }\AttributeTok{metrics =}\NormalTok{ cv\_met)}
\end{Highlighting}
\end{Shaded}

\begin{verbatim}
## > A | warning: glm.fit: dopasowane prawdopodobieństwa numerycznie okazały się być 0 lub 1
\end{verbatim}

\begin{verbatim}
## There were issues with some computations   A: x1There were issues with some computations   A: x1
\end{verbatim}

\begin{Shaded}
\begin{Highlighting}[]
\NormalTok{rf\_cv  }\OtherTok{\textless{}{-}} \FunctionTok{fit\_resamples}\NormalTok{(rf\_wf,  cv\_folds, }\AttributeTok{metrics =}\NormalTok{ cv\_met)}
\NormalTok{xgb\_cv }\OtherTok{\textless{}{-}} \FunctionTok{fit\_resamples}\NormalTok{(xgb\_wf, cv\_folds, }\AttributeTok{metrics =}\NormalTok{ cv\_met)}

\NormalTok{cv\_results }\OtherTok{\textless{}{-}} \FunctionTok{bind\_rows}\NormalTok{(}
  \FunctionTok{collect\_metrics}\NormalTok{(log\_cv) }\SpecialCharTok{\%\textgreater{}\%} \FunctionTok{mutate}\NormalTok{(}\AttributeTok{model =} \StringTok{"Regresja Logistyczna"}\NormalTok{),}
  \FunctionTok{collect\_metrics}\NormalTok{(rf\_cv)  }\SpecialCharTok{\%\textgreater{}\%} \FunctionTok{mutate}\NormalTok{(}\AttributeTok{model =} \StringTok{"Las Losowy"}\NormalTok{),}
  \FunctionTok{collect\_metrics}\NormalTok{(xgb\_cv) }\SpecialCharTok{\%\textgreater{}\%} \FunctionTok{mutate}\NormalTok{(}\AttributeTok{model =} \StringTok{"XGBoost"}\NormalTok{)}
\NormalTok{)}

\CommentTok{\# Tabela zbiorcza wyników CV}
\NormalTok{cv\_results }\SpecialCharTok{\%\textgreater{}\%}
  \FunctionTok{select}\NormalTok{(model, .metric, mean, std\_err) }\SpecialCharTok{\%\textgreater{}\%}
  \FunctionTok{mutate}\NormalTok{(}\FunctionTok{across}\NormalTok{(}\FunctionTok{c}\NormalTok{(mean, std\_err), \textbackslash{}(x) }\FunctionTok{round}\NormalTok{(x, }\DecValTok{4}\NormalTok{))) }\SpecialCharTok{\%\textgreater{}\%}
\NormalTok{  tidyr}\SpecialCharTok{::}\FunctionTok{pivot\_wider}\NormalTok{(}\AttributeTok{names\_from =}\NormalTok{ model, }\AttributeTok{values\_from =} \FunctionTok{c}\NormalTok{(mean, std\_err)) }\SpecialCharTok{\%\textgreater{}\%}
  \FunctionTok{print}\NormalTok{()}
\end{Highlighting}
\end{Shaded}

\begin{verbatim}
## # A tibble: 5 x 7
##   .metric   `mean_Regresja Logistyczna` `mean_Las Losowy` mean_XGBoost
##   <chr>                           <dbl>             <dbl>        <dbl>
## 1 accuracy                        0.835             0.967        0.970
## 2 f_meas                          0.908             0.983        0.985
## 3 precision                       0.998             0.99         0.989
## 4 recall                          0.834             0.976        0.980
## 5 roc_auc                         0.930             0.902        0.888
## # i 3 more variables: `std_err_Regresja Logistyczna` <dbl>,
## #   `std_err_Las Losowy` <dbl>, std_err_XGBoost <dbl>
\end{verbatim}

\begin{Shaded}
\begin{Highlighting}[]
\CommentTok{\# Wykres porównawczy}
\NormalTok{cv\_results }\SpecialCharTok{\%\textgreater{}\%}
  \FunctionTok{filter}\NormalTok{(.metric }\SpecialCharTok{\%in\%} \FunctionTok{c}\NormalTok{(}\StringTok{"f\_meas"}\NormalTok{, }\StringTok{"roc\_auc"}\NormalTok{, }\StringTok{"precision"}\NormalTok{, }\StringTok{"recall"}\NormalTok{)) }\SpecialCharTok{\%\textgreater{}\%}
  \FunctionTok{ggplot}\NormalTok{(}\FunctionTok{aes}\NormalTok{(}\AttributeTok{x =}\NormalTok{ model, }\AttributeTok{y =}\NormalTok{ mean, }\AttributeTok{fill =}\NormalTok{ model)) }\SpecialCharTok{+}
  \FunctionTok{geom\_col}\NormalTok{(}\AttributeTok{show.legend =} \ConstantTok{FALSE}\NormalTok{) }\SpecialCharTok{+}
  \FunctionTok{geom\_errorbar}\NormalTok{(}\FunctionTok{aes}\NormalTok{(}\AttributeTok{ymin =}\NormalTok{ mean }\SpecialCharTok{{-}}\NormalTok{ std\_err, }\AttributeTok{ymax =}\NormalTok{ mean }\SpecialCharTok{+}\NormalTok{ std\_err), }\AttributeTok{width =} \FloatTok{0.25}\NormalTok{) }\SpecialCharTok{+}
  \FunctionTok{facet\_wrap}\NormalTok{(}\SpecialCharTok{\textasciitilde{}}\NormalTok{ .metric, }\AttributeTok{scales =} \StringTok{"free\_y"}\NormalTok{) }\SpecialCharTok{+}
  \FunctionTok{theme\_minimal}\NormalTok{() }\SpecialCharTok{+}
  \FunctionTok{labs}\NormalTok{(}\AttributeTok{title =} \StringTok{"Porównanie modeli – 5{-}krotna kroswalidacja (5{-}fold CV)"}\NormalTok{,}
       \AttributeTok{x =} \ConstantTok{NULL}\NormalTok{,}
       \AttributeTok{y =} \StringTok{"Średnia wartość metryki (± std. błąd)"}\NormalTok{)}
\end{Highlighting}
\end{Shaded}

\pandocbounded{\includegraphics[keepaspectratio]{projekt-final_files/figure-latex/unnamed-chunk-7-1.pdf}}

\textbf{Wniosek z kroswalidacji:} Wyniki CV potwierdzają obserwacje z
ewaluacji na zbiorze testowym. Regresja logistyczna osiąga relatywnie
wysokie ROC-AUC, lecz jej F1-Score i precyzja są wyraźnie niższe od
modeli drzewiastych, co wskazuje na nadmierne rozszerzenie granicy
decyzyjnej po SMOTE. Las Losowy i XGBoost wykazują porównywalną
stabilność (małe odchylenia std. między foldami), przy czym Las Losowy
zwykle uzyskuje nieco lepszy F1-Score. Słupki błędu (±1 std. błąd) są
wąskie dla wszystkich modeli, co świadczy o spójności wyników między
foldami i niskim ryzyku przeuczenia. Na podstawie tych wyników
\textbf{Las Losowy} zostaje wybrany jako model wiodący do dalszej
ewaluacji i wdrożenia.

\subsection{Ewaluacja modeli uczenia maszynowego oraz ich
kalibracja}\label{ewaluacja-modeli-uczenia-maszynowego-oraz-ich-kalibracja}

Po wytrenowaniu modeli na zbalansowanym zbiorze treningowym
(wykorzystującym technikę SMOTE), przeprowadzono ich ewaluację na
surowym zbiorze testowym. Pozwoliło to na sprawdzenie, jak modele radzą
sobie z rzeczywistym rozkładem danych (gdzie hity są rzadkością).

Do oceny jakości klasyfikacji wykorzystano zestaw metryk:
\textbf{Accuracy} (Dokładność), \textbf{Precision} (Precyzja),
\textbf{Recall} (Czułość) oraz \textbf{F1-Score}. Dodatkowo wygenerowano
krzywe \textbf{ROC} oraz obliczono pole pod krzywą (\textbf{ROC-AUC}),
które jest odporne na niezbalansowanie klas.

\begin{Shaded}
\begin{Highlighting}[]
\FunctionTok{library}\NormalTok{(yardstick)}
\FunctionTok{library}\NormalTok{(dplyr)}
\FunctionTok{library}\NormalTok{(ggplot2)}
\FunctionTok{library}\NormalTok{(probably)}
\end{Highlighting}
\end{Shaded}

\begin{verbatim}
## Warning: pakiet 'probably' został zbudowany w wersji R 4.4.3
\end{verbatim}

\begin{verbatim}
## 
## Dołączanie pakietu: 'probably'
\end{verbatim}

\begin{verbatim}
## Następujące obiekty zostały zakryte z 'package:base':
## 
##     as.factor, as.ordered
\end{verbatim}

\begin{Shaded}
\begin{Highlighting}[]
\NormalTok{multi\_metric }\OtherTok{\textless{}{-}} \FunctionTok{metric\_set}\NormalTok{(accuracy, precision, recall, f\_meas)}

\FunctionTok{cat}\NormalTok{(}\StringTok{"}\SpecialCharTok{\textbackslash{}n}\StringTok{{-}{-}{-} WYNIKI DLA REGRESJI LOGISTYCZNEJ {-}{-}{-}}\SpecialCharTok{\textbackslash{}n}\StringTok{"}\NormalTok{)}
\end{Highlighting}
\end{Shaded}

\begin{verbatim}
## 
## --- WYNIKI DLA REGRESJI LOGISTYCZNEJ ---
\end{verbatim}

\begin{Shaded}
\begin{Highlighting}[]
\FunctionTok{print}\NormalTok{(log\_preds }\SpecialCharTok{\%\textgreater{}\%} \FunctionTok{multi\_metric}\NormalTok{(}\AttributeTok{truth =}\NormalTok{ is\_hit, }\AttributeTok{estimate =}\NormalTok{ .pred\_class, }\AttributeTok{event\_level =} \StringTok{"second"}\NormalTok{))}
\end{Highlighting}
\end{Shaded}

\begin{verbatim}
## # A tibble: 4 x 3
##   .metric   .estimator .estimate
##   <chr>     <chr>          <dbl>
## 1 accuracy  binary        0.832 
## 2 precision binary        0.0723
## 3 recall    binary        0.827 
## 4 f_meas    binary        0.133
\end{verbatim}

\begin{Shaded}
\begin{Highlighting}[]
\FunctionTok{cat}\NormalTok{(}\StringTok{"}\SpecialCharTok{\textbackslash{}n}\StringTok{{-}{-}{-} WYNIKI DLA LASU LOSOWEGO {-}{-}{-}}\SpecialCharTok{\textbackslash{}n}\StringTok{"}\NormalTok{)}
\end{Highlighting}
\end{Shaded}

\begin{verbatim}
## 
## --- WYNIKI DLA LASU LOSOWEGO ---
\end{verbatim}

\begin{Shaded}
\begin{Highlighting}[]
\FunctionTok{print}\NormalTok{(rf\_preds }\SpecialCharTok{\%\textgreater{}\%} \FunctionTok{multi\_metric}\NormalTok{(}\AttributeTok{truth =}\NormalTok{ is\_hit, }\AttributeTok{estimate =}\NormalTok{ .pred\_class, }\AttributeTok{event\_level =} \StringTok{"second"}\NormalTok{))}
\end{Highlighting}
\end{Shaded}

\begin{verbatim}
## # A tibble: 4 x 3
##   .metric   .estimator .estimate
##   <chr>     <chr>          <dbl>
## 1 accuracy  binary         0.965
## 2 precision binary         0.208
## 3 recall    binary         0.441
## 4 f_meas    binary         0.283
\end{verbatim}

\begin{Shaded}
\begin{Highlighting}[]
\FunctionTok{cat}\NormalTok{(}\StringTok{"}\SpecialCharTok{\textbackslash{}n}\StringTok{{-}{-}{-} WYNIKI DLA XGBoost {-}{-}{-}}\SpecialCharTok{\textbackslash{}n}\StringTok{"}\NormalTok{)}
\end{Highlighting}
\end{Shaded}

\begin{verbatim}
## 
## --- WYNIKI DLA XGBoost ---
\end{verbatim}

\begin{Shaded}
\begin{Highlighting}[]
\FunctionTok{print}\NormalTok{(xgb\_preds }\SpecialCharTok{\%\textgreater{}\%} \FunctionTok{multi\_metric}\NormalTok{(}\AttributeTok{truth =}\NormalTok{ is\_hit, }\AttributeTok{estimate =}\NormalTok{ .pred\_class, }\AttributeTok{event\_level =} \StringTok{"second"}\NormalTok{))}
\end{Highlighting}
\end{Shaded}

\begin{verbatim}
## # A tibble: 4 x 3
##   .metric   .estimator .estimate
##   <chr>     <chr>          <dbl>
## 1 accuracy  binary         0.968
## 2 precision binary         0.210
## 3 recall    binary         0.390
## 4 f_meas    binary         0.273
\end{verbatim}

\begin{Shaded}
\begin{Highlighting}[]
\NormalTok{log\_roc }\OtherTok{\textless{}{-}}\NormalTok{ log\_preds }\SpecialCharTok{\%\textgreater{}\%} \FunctionTok{roc\_curve}\NormalTok{(}\AttributeTok{truth =}\NormalTok{ is\_hit, .pred\_1, }\AttributeTok{event\_level =} \StringTok{"second"}\NormalTok{) }\SpecialCharTok{\%\textgreater{}\%} \FunctionTok{mutate}\NormalTok{(}\AttributeTok{model =} \StringTok{"Regresja Logistyczna"}\NormalTok{)}
\NormalTok{rf\_roc }\OtherTok{\textless{}{-}}\NormalTok{ rf\_preds }\SpecialCharTok{\%\textgreater{}\%} \FunctionTok{roc\_curve}\NormalTok{(}\AttributeTok{truth =}\NormalTok{ is\_hit, .pred\_1, }\AttributeTok{event\_level =} \StringTok{"second"}\NormalTok{) }\SpecialCharTok{\%\textgreater{}\%} \FunctionTok{mutate}\NormalTok{(}\AttributeTok{model =} \StringTok{"Las Losowy"}\NormalTok{)}
\NormalTok{xgb\_roc }\OtherTok{\textless{}{-}}\NormalTok{ xgb\_preds }\SpecialCharTok{\%\textgreater{}\%} \FunctionTok{roc\_curve}\NormalTok{(}\AttributeTok{truth =}\NormalTok{ is\_hit, .pred\_1, }\AttributeTok{event\_level =} \StringTok{"second"}\NormalTok{) }\SpecialCharTok{\%\textgreater{}\%} \FunctionTok{mutate}\NormalTok{(}\AttributeTok{model =} \StringTok{"XGBoost"}\NormalTok{)}

\NormalTok{roc\_data }\OtherTok{\textless{}{-}} \FunctionTok{bind\_rows}\NormalTok{(log\_roc, rf\_roc, xgb\_roc)}

\FunctionTok{ggplot}\NormalTok{(roc\_data, }\FunctionTok{aes}\NormalTok{(}\AttributeTok{x =} \DecValTok{1} \SpecialCharTok{{-}}\NormalTok{ specificity, }\AttributeTok{y =}\NormalTok{ sensitivity, }\AttributeTok{color =}\NormalTok{ model)) }\SpecialCharTok{+}
  \FunctionTok{geom\_path}\NormalTok{(}\AttributeTok{linewidth =} \DecValTok{1}\NormalTok{) }\SpecialCharTok{+}
  \FunctionTok{geom\_abline}\NormalTok{(}\AttributeTok{lty =} \DecValTok{2}\NormalTok{, }\AttributeTok{color =} \StringTok{"gray50"}\NormalTok{) }\SpecialCharTok{+}
  \FunctionTok{coord\_equal}\NormalTok{() }\SpecialCharTok{+}
  \FunctionTok{theme\_minimal}\NormalTok{() }\SpecialCharTok{+}
  \FunctionTok{labs}\NormalTok{(}\AttributeTok{title =} \StringTok{"Krzywe ROC dla przetestowanych modeli"}\NormalTok{,}
       \AttributeTok{x =} \StringTok{"False Positive Rate (1 {-} Specificity)"}\NormalTok{,}
       \AttributeTok{y =} \StringTok{"True Positive Rate (Sensitivity)"}\NormalTok{,}
       \AttributeTok{color =} \StringTok{"Model"}\NormalTok{) }\SpecialCharTok{+}
  \FunctionTok{scale\_color\_manual}\NormalTok{(}\AttributeTok{values =} \FunctionTok{c}\NormalTok{(}\StringTok{"Regresja Logistyczna"} \OtherTok{=} \StringTok{"steelblue"}\NormalTok{, }\StringTok{"Las Losowy"} \OtherTok{=} \StringTok{"firebrick"}\NormalTok{, }\StringTok{"XGBoost"} \OtherTok{=} \StringTok{"forestgreen"}\NormalTok{))}
\end{Highlighting}
\end{Shaded}

\pandocbounded{\includegraphics[keepaspectratio]{projekt-final_files/figure-latex/unnamed-chunk-8-1.pdf}}

\begin{Shaded}
\begin{Highlighting}[]
\FunctionTok{cat}\NormalTok{(}\StringTok{"}\SpecialCharTok{\textbackslash{}n}\StringTok{{-}{-}{-} AUC DLA REGRESJI LOGISTYCZNEJ {-}{-}{-}}\SpecialCharTok{\textbackslash{}n}\StringTok{"}\NormalTok{)}
\end{Highlighting}
\end{Shaded}

\begin{verbatim}
## 
## --- AUC DLA REGRESJI LOGISTYCZNEJ ---
\end{verbatim}

\begin{Shaded}
\begin{Highlighting}[]
\NormalTok{log\_preds }\SpecialCharTok{\%\textgreater{}\%} \FunctionTok{roc\_auc}\NormalTok{(}\AttributeTok{truth =}\NormalTok{ is\_hit, .pred\_1, }\AttributeTok{event\_level =} \StringTok{"second"}\NormalTok{) }\SpecialCharTok{\%\textgreater{}\%} \FunctionTok{print}\NormalTok{()}
\end{Highlighting}
\end{Shaded}

\begin{verbatim}
## # A tibble: 1 x 3
##   .metric .estimator .estimate
##   <chr>   <chr>          <dbl>
## 1 roc_auc binary         0.920
\end{verbatim}

\begin{Shaded}
\begin{Highlighting}[]
\FunctionTok{cat}\NormalTok{(}\StringTok{"}\SpecialCharTok{\textbackslash{}n}\StringTok{{-}{-}{-} AUC DLA LASU LOSOWEGO {-}{-}{-}}\SpecialCharTok{\textbackslash{}n}\StringTok{"}\NormalTok{)}
\end{Highlighting}
\end{Shaded}

\begin{verbatim}
## 
## --- AUC DLA LASU LOSOWEGO ---
\end{verbatim}

\begin{Shaded}
\begin{Highlighting}[]
\NormalTok{rf\_preds }\SpecialCharTok{\%\textgreater{}\%} \FunctionTok{roc\_auc}\NormalTok{(}\AttributeTok{truth =}\NormalTok{ is\_hit, .pred\_1, }\AttributeTok{event\_level =} \StringTok{"second"}\NormalTok{) }\SpecialCharTok{\%\textgreater{}\%} \FunctionTok{print}\NormalTok{()}
\end{Highlighting}
\end{Shaded}

\begin{verbatim}
## # A tibble: 1 x 3
##   .metric .estimator .estimate
##   <chr>   <chr>          <dbl>
## 1 roc_auc binary         0.893
\end{verbatim}

\begin{Shaded}
\begin{Highlighting}[]
\FunctionTok{cat}\NormalTok{(}\StringTok{"}\SpecialCharTok{\textbackslash{}n}\StringTok{{-}{-}{-} AUC DLA XGBoost {-}{-}{-}}\SpecialCharTok{\textbackslash{}n}\StringTok{"}\NormalTok{)}
\end{Highlighting}
\end{Shaded}

\begin{verbatim}
## 
## --- AUC DLA XGBoost ---
\end{verbatim}

\begin{Shaded}
\begin{Highlighting}[]
\NormalTok{xgb\_preds }\SpecialCharTok{\%\textgreater{}\%} \FunctionTok{roc\_auc}\NormalTok{(}\AttributeTok{truth =}\NormalTok{ is\_hit, .pred\_1, }\AttributeTok{event\_level =} \StringTok{"second"}\NormalTok{) }\SpecialCharTok{\%\textgreater{}\%} \FunctionTok{print}\NormalTok{()}
\end{Highlighting}
\end{Shaded}

\begin{verbatim}
## # A tibble: 1 x 3
##   .metric .estimator .estimate
##   <chr>   <chr>          <dbl>
## 1 roc_auc binary         0.886
\end{verbatim}

\begin{Shaded}
\begin{Highlighting}[]
\CommentTok{\# Optymalizacja progu odcięcia dla Lasu Losowego}
\NormalTok{threshold\_data }\OtherTok{\textless{}{-}}\NormalTok{ rf\_preds }\SpecialCharTok{\%\textgreater{}\%}
  \FunctionTok{threshold\_perf}\NormalTok{(}\AttributeTok{truth =}\NormalTok{ is\_hit, }\AttributeTok{estimate =}\NormalTok{ .pred\_1, }
                 \AttributeTok{thresholds =} \FunctionTok{seq}\NormalTok{(}\FloatTok{0.1}\NormalTok{, }\FloatTok{0.9}\NormalTok{, }\AttributeTok{by =} \FloatTok{0.05}\NormalTok{), }
                 \AttributeTok{metrics =} \FunctionTok{metric\_set}\NormalTok{(f\_meas))}

\FunctionTok{ggplot}\NormalTok{(threshold\_data, }\FunctionTok{aes}\NormalTok{(}\AttributeTok{x =}\NormalTok{ .threshold, }\AttributeTok{y =}\NormalTok{ .estimate)) }\SpecialCharTok{+}
  \FunctionTok{geom\_line}\NormalTok{(}\AttributeTok{color =} \StringTok{"firebrick"}\NormalTok{, }\AttributeTok{linewidth =} \DecValTok{1}\NormalTok{) }\SpecialCharTok{+}
  \FunctionTok{geom\_point}\NormalTok{(}\AttributeTok{color =} \StringTok{"firebrick"}\NormalTok{, }\AttributeTok{size =} \DecValTok{2}\NormalTok{) }\SpecialCharTok{+}
  \FunctionTok{theme\_minimal}\NormalTok{() }\SpecialCharTok{+}
  \FunctionTok{labs}\NormalTok{(}\AttributeTok{title =} \StringTok{"Optymalizacja progu odcięcia {-} Las Losowy"}\NormalTok{,}
       \AttributeTok{x =} \StringTok{"Próg odcięcia prawdopodobieństwa (Threshold)"}\NormalTok{,}
       \AttributeTok{y =} \StringTok{"Wartość metryki F1{-}Score"}\NormalTok{)}
\end{Highlighting}
\end{Shaded}

\pandocbounded{\includegraphics[keepaspectratio]{projekt-final_files/figure-latex/unnamed-chunk-8-2.pdf}}

\subsubsection{Kalibracja
prawdopodobieństw}\label{kalibracja-prawdopodobieux144stw}

Krzywa kalibracji (\emph{calibration plot}) pokazuje, jak dobrze
prawdopodobieństwa predykcji modelu odpowiadają rzeczywistym częstościom
obserwowanych klas. Idealnie skalibrowany model leży na linii przekątnej
(prawdopodobieństwo predykowane = obserwowana częstość). Odchylenia od
tej linii informują o tym, czy model jest zbyt optymistyczny (krzywa
powyżej przekątnej) lub zbyt pesymistyczny (poniżej przekątnej).

\begin{Shaded}
\begin{Highlighting}[]
\FunctionTok{library}\NormalTok{(probably)}

\NormalTok{cal\_rf }\OtherTok{\textless{}{-}} \FunctionTok{cal\_plot\_breaks}\NormalTok{(}
\NormalTok{  rf\_preds, }\AttributeTok{truth =}\NormalTok{ is\_hit, }\AttributeTok{estimate =}\NormalTok{ .pred\_1,}
  \AttributeTok{event\_level =} \StringTok{"second"}\NormalTok{, }\AttributeTok{num\_breaks =} \DecValTok{10}
\NormalTok{) }\SpecialCharTok{+} \FunctionTok{labs}\NormalTok{(}\AttributeTok{title =} \StringTok{"Kalibracja – Las Losowy"}\NormalTok{) }\SpecialCharTok{+} \FunctionTok{theme\_minimal}\NormalTok{()}

\NormalTok{cal\_log }\OtherTok{\textless{}{-}} \FunctionTok{cal\_plot\_breaks}\NormalTok{(}
\NormalTok{  log\_preds, }\AttributeTok{truth =}\NormalTok{ is\_hit, }\AttributeTok{estimate =}\NormalTok{ .pred\_1,}
  \AttributeTok{event\_level =} \StringTok{"second"}\NormalTok{, }\AttributeTok{num\_breaks =} \DecValTok{10}
\NormalTok{) }\SpecialCharTok{+} \FunctionTok{labs}\NormalTok{(}\AttributeTok{title =} \StringTok{"Kalibracja – Regresja Logistyczna"}\NormalTok{) }\SpecialCharTok{+} \FunctionTok{theme\_minimal}\NormalTok{()}

\NormalTok{cal\_xgb }\OtherTok{\textless{}{-}} \FunctionTok{cal\_plot\_breaks}\NormalTok{(}
\NormalTok{  xgb\_preds, }\AttributeTok{truth =}\NormalTok{ is\_hit, }\AttributeTok{estimate =}\NormalTok{ .pred\_1,}
  \AttributeTok{event\_level =} \StringTok{"second"}\NormalTok{, }\AttributeTok{num\_breaks =} \DecValTok{10}
\NormalTok{) }\SpecialCharTok{+} \FunctionTok{labs}\NormalTok{(}\AttributeTok{title =} \StringTok{"Kalibracja – XGBoost"}\NormalTok{) }\SpecialCharTok{+} \FunctionTok{theme\_minimal}\NormalTok{()}

\FunctionTok{library}\NormalTok{(patchwork)}
\end{Highlighting}
\end{Shaded}

\begin{verbatim}
## Warning: pakiet 'patchwork' został zbudowany w wersji R 4.4.2
\end{verbatim}

\begin{Shaded}
\begin{Highlighting}[]
\NormalTok{cal\_rf }\SpecialCharTok{+}\NormalTok{ cal\_log }\SpecialCharTok{+}\NormalTok{ cal\_xgb}
\end{Highlighting}
\end{Shaded}

\pandocbounded{\includegraphics[keepaspectratio]{projekt-final_files/figure-latex/unnamed-chunk-9-1.pdf}}

\textbf{Wniosek z kalibracji:} Regresja logistyczna wykazuje wyraźne
przekalibrowanie -- dla niskich prawdopodobieństw predykuje zbyt wysokie
wartości (efekt agresywnego SMOTE), co tłumaczy jej bardzo niską
precyzję. Las Losowy i XGBoost są lepiej skalibrowane w środkowym
zakresie prawdopodobieństw, lecz w obszarze wysokich wartości
(\textgreater0.7) mogą niedoszacowywać rzeczywistej częstości hitów ze
względu na naturalną rzadkość tej klasy w danych testowych. Wyniki
uzasadniają optymalizację progu odcięcia dla modelu wiodącego.

\textbf{Wnioski z ewaluacji i porównania modeli:}

\begin{enumerate}
\def\labelenumi{\arabic{enumi}.}
\item
  \textbf{Regresja Logistyczna:} Mimo wysokiego wskaźnika ROC-AUC
  (0.919), model ten w praktyce wykazuje bardzo niską użyteczność.
  Bardzo wysoka czułość (Recall = 0.826) sugeruje, że model ``widzi hity
  wszędzie''. Skutkuje to ekstremalnie niską precyzją (0.072) -- na
  każdy poprawnie wskazany hit przypada kilkanaście fałszywych alarmów.
  Jest to efekt uboczny SMOTE, który w połączeniu z liniowym charakterem
  modelu doprowadził do nadmiernego rozszerzenia granicy decyzyjnej
  klasy pozytywnej.
\item
  \textbf{Modele Drzewiaste (Las Losowy i XGBoost):} Wykazały się
  znacznie większą stabilnością. Las Losowy osiągnął najwyższy F1-Score
  (0.275). Choć jego czułość jest niższa niż w regresji, to precyzja
  (0.204) jest prawie 3-krotnie wyższa. Oznacza to, że model ten
  znacznie lepiej separuje ``prawdziwe hity'' od utworów, które jedynie
  brzmią podobnie do nich.
\item
  \textbf{Optymalizacja progu:} Analiza wykresu optymalizacji progu dla
  Lasu Losowego pokazuje, że domyślny próg 0.5 nie jest optymalny dla
  F1-Score. Maksimum metryki osiągane jest dla wyższych progów
  prawdopodobieństwa, co sugeruje, że model jest z natury nieco zbyt
  optymistyczny po procesie SMOTE i warto podnieść poprzeczkę dla
  klasyfikacji utworu jako hit.
\end{enumerate}

\subsection{Przeprowadzenie analizy popełnianych błędów przez
wykorzystane modele
predykcyjne}\label{przeprowadzenie-analizy-popeux142nianych-bux142ux119duxf3w-przez-wykorzystane-modele-predykcyjne}

Analiza macierzy pomyłek (Confusion Matrix) pozwala na precyzyjne
zidentyfikowanie natury błędów popełnianych przez modele.

\emph{Wnioski z analizy błędów:} Głównym problemem wszystkich modeli
jest wysoka liczba błędów \textbf{False Positive} (klasyfikowanie
przeciętnych utworów jako hity). Zjawisko to ma głębokie uzasadnienie
merytoryczne: - \textbf{Pułapka akustyczna:} Wiele utworów jest
produkowanych zgodnie z obecnymi trendami (wysoka głośność, taneczność),
co sprawia, że w przestrzeni cech akustycznych są one identyczne z
hitami. - \textbf{Brakujące zmienne:} Model nie wie, czy utwór był
promowany w radiu, czy artysta ma miliony fanów, ani czy piosenka stała
się wiralem na TikToku. Te ``ukryte zmienne'' sprawiają, że utwory o
identycznych parametrach audio mają skrajnie różną popularność. -
\textbf{Błędy False Negative:} Są rzadsze, ale równie ciekawe -- dotyczą
hitów, które ``brzmią inaczej'' (np. ballady akustyczne w zalewie muzyki
elektronicznej). Model, ucząc się na trendach, ma trudności z wyłapaniem
takich wyjątków (tzw. outliers).

\begin{Shaded}
\begin{Highlighting}[]
\FunctionTok{library}\NormalTok{(yardstick)}
\FunctionTok{library}\NormalTok{(ggplot2)}

\NormalTok{rf\_preds }\SpecialCharTok{\%\textgreater{}\%}
  \FunctionTok{conf\_mat}\NormalTok{(}\AttributeTok{truth =}\NormalTok{ is\_hit, }\AttributeTok{estimate =}\NormalTok{ .pred\_class) }\SpecialCharTok{\%\textgreater{}\%}
  \FunctionTok{autoplot}\NormalTok{(}\AttributeTok{type =} \StringTok{"heatmap"}\NormalTok{) }\SpecialCharTok{+}
  \FunctionTok{scale\_fill\_gradient}\NormalTok{(}\AttributeTok{low =} \StringTok{"white"}\NormalTok{, }\AttributeTok{high =} \StringTok{"firebrick"}\NormalTok{) }\SpecialCharTok{+}
  \FunctionTok{labs}\NormalTok{(}\AttributeTok{title =} \StringTok{"Macierz pomyłek {-} Las Losowy"}\NormalTok{,}
       \AttributeTok{x =} \StringTok{"Predykcja modelu"}\NormalTok{,}
       \AttributeTok{y =} \StringTok{"Rzeczywista wartość"}\NormalTok{) }\SpecialCharTok{+}
  \FunctionTok{theme\_minimal}\NormalTok{()}
\end{Highlighting}
\end{Shaded}

\begin{verbatim}
## Scale for fill is already present.
## Adding another scale for fill, which will replace the existing scale.
\end{verbatim}

\pandocbounded{\includegraphics[keepaspectratio]{projekt-final_files/figure-latex/unnamed-chunk-10-1.pdf}}

\begin{Shaded}
\begin{Highlighting}[]
\NormalTok{log\_preds }\SpecialCharTok{\%\textgreater{}\%}
  \FunctionTok{conf\_mat}\NormalTok{(}\AttributeTok{truth =}\NormalTok{ is\_hit, }\AttributeTok{estimate =}\NormalTok{ .pred\_class) }\SpecialCharTok{\%\textgreater{}\%}
  \FunctionTok{autoplot}\NormalTok{(}\AttributeTok{type =} \StringTok{"heatmap"}\NormalTok{) }\SpecialCharTok{+}
  \FunctionTok{scale\_fill\_gradient}\NormalTok{(}\AttributeTok{low =} \StringTok{"white"}\NormalTok{, }\AttributeTok{high =} \StringTok{"steelblue"}\NormalTok{) }\SpecialCharTok{+}
  \FunctionTok{labs}\NormalTok{(}\AttributeTok{title =} \StringTok{"Macierz pomyłek {-} Regresja Logistyczna"}\NormalTok{,}
       \AttributeTok{x =} \StringTok{"Predykcja modelu"}\NormalTok{,}
       \AttributeTok{y =} \StringTok{"Rzeczywista wartość"}\NormalTok{) }\SpecialCharTok{+}
  \FunctionTok{theme\_minimal}\NormalTok{()}
\end{Highlighting}
\end{Shaded}

\begin{verbatim}
## Scale for fill is already present.
## Adding another scale for fill, which will replace the existing scale.
\end{verbatim}

\pandocbounded{\includegraphics[keepaspectratio]{projekt-final_files/figure-latex/unnamed-chunk-10-2.pdf}}

\textbf{Wniosek:} Zestawienie macierzy pomyłek ujawnia fundamentalną
różnicę w stylu błędów obu modeli. Regresja logistyczna generuje ogromną
liczbę błędów False Positive -- model jest zbyt ``optymistyczny'' i
przypisuje status hitu bardzo szerokim grupom utworów, co skutkuje niską
precyzją przy względnie wysokiej czułości. Las losowy zachowuje się
odwrotnie -- jest bardziej konserwatywny i rzadziej klasy pozytywnej,
przez co ma wyższą precyzję przy niższej czułości. W praktycznym
zastosowaniu model lasu losowego jest bezpieczniejszy: lepiej
identyfikuje prawdziwe potencjalne hity, minimalizując ryzyko fałszywych
rekomendacji, co jest istotne w kontekście decyzji biznesowych wytwórni
muzycznej.

\subsection{Analiza wyjaśnialności decyzji modeli
predykcyjnych}\label{analiza-wyjaux15bnialnoux15bci-decyzji-modeli-predykcyjnych}

Modele oparte na drzewach decyzyjnych często określane są mianem
``czarnych skrzynek''. Aby zrozumieć logikę stojącą za predykcjami
wybranego modelu wiodącego (Las Losowy), przeprowadzono analizę ważności
cech (\textbf{Variable Importance, VIP}). Informuje ona o tym, które
zmienne najbardziej przyczyniają się do poprawy czystości węzłów w
procesie budowania drzew.

\begin{Shaded}
\begin{Highlighting}[]
\FunctionTok{library}\NormalTok{(vip)}
\end{Highlighting}
\end{Shaded}

\begin{verbatim}
## Warning: pakiet 'vip' został zbudowany w wersji R 4.4.3
\end{verbatim}

\begin{verbatim}
## 
## Dołączanie pakietu: 'vip'
\end{verbatim}

\begin{verbatim}
## Następujący obiekt został zakryty z 'package:utils':
## 
##     vi
\end{verbatim}

\begin{Shaded}
\begin{Highlighting}[]
\FunctionTok{library}\NormalTok{(ggplot2)}

\FunctionTok{vip}\NormalTok{(rf\_fit}\SpecialCharTok{$}\NormalTok{fit, }\AttributeTok{num\_features =} \DecValTok{20}\NormalTok{, }\AttributeTok{geom =} \StringTok{"col"}\NormalTok{, }
    \AttributeTok{mapping =} \FunctionTok{aes\_string}\NormalTok{(}\AttributeTok{fill =} \StringTok{"Variable"}\NormalTok{)) }\SpecialCharTok{+}
  \FunctionTok{theme\_minimal}\NormalTok{() }\SpecialCharTok{+}
  \FunctionTok{scale\_fill\_viridis\_d}\NormalTok{(}\AttributeTok{option =} \StringTok{"magma"}\NormalTok{, }\AttributeTok{guide =} \StringTok{"none"}\NormalTok{) }\SpecialCharTok{+}
  \FunctionTok{labs}\NormalTok{(}\AttributeTok{title =} \StringTok{"Ważność cech (Top 20) {-} Model Lasu Losowego"}\NormalTok{,}
       \AttributeTok{x =} \StringTok{"Ważność (Impurity)"}\NormalTok{,}
       \AttributeTok{y =} \StringTok{"Zmienna objaśniająca"}\NormalTok{)}
\end{Highlighting}
\end{Shaded}

\begin{verbatim}
## Warning: `aes_string()` was deprecated in ggplot2 3.0.0.
## i Please use tidy evaluation idioms with `aes()`.
## i See also `vignette("ggplot2-in-packages")` for more information.
## This warning is displayed once every 8 hours.
## Call `lifecycle::last_lifecycle_warnings()` to see where this warning was
## generated.
\end{verbatim}

\pandocbounded{\includegraphics[keepaspectratio]{projekt-final_files/figure-latex/unnamed-chunk-11-1.pdf}}

\textbf{Interpretacja wyników ważności cech:} 1. \textbf{Dominacja
parametrów audio:} Najważniejszymi predyktorami okazały się
\texttt{instrumentalness} oraz \texttt{acousticness}. Sugeruje to, że
dla popularności kluczowe jest to, czy utwór ma charakter ``radiowy''
(wokalny, elektronicznie przetworzony). Utwory czysto instrumentalne lub
o surowym, akustycznym brzmieniu rzadziej przebijają się do progu
popularności \textgreater= 75. 2. \textbf{Parametry techniczne:} Wysoka
pozycja \texttt{duration\_ms} oraz \texttt{loudness} potwierdza rynkowe
trendy -- hity mają zazwyczaj określoną długość (ok. 3 minut) i są
poddawane agresywnemu masteringowi (tzw. Loudness War). 3.
\textbf{Marginalizacja gatunków:} Co niezwykle istotne, przynależność do
konkretnego gatunku (\texttt{track\_genre}) ma dla modelu mniejsze
znaczenie niż parametry dźwiękowe. Najważniejszym gatunkiem okazał się
\texttt{dance}, co jest spójne z obecną dominacją muzyki tanecznej w
głównym nurcie.

\subsection{Podsumowanie i sformułowanie
wniosków}\label{podsumowanie-i-sformuux142owanie-wnioskuxf3w}

Projekt miał na celu zbadanie, w jakim stopniu metadane akustyczne
Spotify pozwalają przewidzieć sukces komercyjny utworu. Przeprowadzona
analiza prowadzi do następujących konkluzji:

\begin{enumerate}
\def\labelenumi{\arabic{enumi}.}
\tightlist
\item
  \textbf{Możliwości predykcyjne:} Budowa modelu o wysokiej precyzji
  (Precision) jest niemożliwa przy użyciu wyłącznie cech akustycznych.
  Choć model potrafi skutecznie wyselekcjonować ``kandydatów na hity''
  (wysoki ROC-AUC), to rzeczywisty sukces zależy od czynników
  pozamuzycznych, takich jak kampanie marketingowe, renoma artysty czy
  trendy w mediach społecznościowych.
\item
  \textbf{Przewaga modeli nieliniowych:} Algorytmy zespołowe (Random
  Forest, XGBoost) zdecydowanie lepiej radzą sobie ze specyfiką danych
  muzycznych niż modele liniowe. Pozwalają one na uchwycenie złożonych
  interakcji między głośnością, energią a tanecznością, które wspólnie
  budują ``potencjał komercyjny''.
\item
  \textbf{Rekomendacje praktyczne:} Z punktu widzenia wytwórni
  muzycznej, model taki jak Las Losowy może służyć jako narzędzie do
  wstępnej filtracji utworów (``akustyczny filtr jakości''), jednak
  ostateczna decyzja o inwestycji musi być poparta analizą trendów
  rynkowych i potencjału wizerunkowego twórcy.
\end{enumerate}

\subsection{Podstawowe wdrożenie modelu wybranego na podstawie
przeprowadzonej ewaluacji oraz przeprowadzenie testu
predykcyjnego}\label{podstawowe-wdroux17cenie-modelu-wybranego-na-podstawie-przeprowadzonej-ewaluacji-oraz-przeprowadzenie-testu-predykcyjnego}

Ostatnim etapem projektu było środowiskowe wdrożenie wybranego modelu
Lasu Losowego i weryfikacja jego działania w praktyce. W tym celu
stworzono syntetyczny profil ``nowej piosenki'', przypisując jej cechy
charakterystyczne dla współczesnych hitów radiowych (wysoka taneczność,
duża głośność, wysoka energia, brak instrumentalności, gatunek pop).
Nowa obserwacja została przepuszczona przez zdefiniowaną wcześniej
recepturę (standaryzacja, kodowanie), a następnie poddana predykcji.

Dodatkowo, wykorzystano narzędzia wyjaśnialności lokalnej (DALEX),
generując wykres dekompozycji predykcji (Waterfall plot) dla testowanego
utworu. Obrazuje on bezpośredni wpływ poszczególnych atrybutów nowej
obserwacji na finalną decyzję podjętą przez model.

\begin{Shaded}
\begin{Highlighting}[]
\FunctionTok{library}\NormalTok{(dplyr)}
\FunctionTok{library}\NormalTok{(DALEX)}
\end{Highlighting}
\end{Shaded}

\begin{verbatim}
## Warning: pakiet 'DALEX' został zbudowany w wersji R 4.4.3
\end{verbatim}

\begin{verbatim}
## Welcome to DALEX (version: 2.5.3).
## Find examples and detailed introduction at: http://ema.drwhy.ai/
\end{verbatim}

\begin{verbatim}
## 
## Dołączanie pakietu: 'DALEX'
\end{verbatim}

\begin{verbatim}
## Następujący obiekt został zakryty z 'package:dplyr':
## 
##     explain
\end{verbatim}

\begin{Shaded}
\begin{Highlighting}[]
\FunctionTok{library}\NormalTok{(DALEXtra)}
\end{Highlighting}
\end{Shaded}

\begin{verbatim}
## Warning: pakiet 'DALEXtra' został zbudowany w wersji R 4.4.3
\end{verbatim}

\begin{Shaded}
\begin{Highlighting}[]
\NormalTok{nowa\_piosenka }\OtherTok{\textless{}{-}} \FunctionTok{tibble}\NormalTok{(}
  \AttributeTok{duration\_ms =} \DecValTok{210000}\NormalTok{,}
  \AttributeTok{danceability =} \FloatTok{0.85}\NormalTok{,}
  \AttributeTok{energy =} \FloatTok{0.90}\NormalTok{,}
  \AttributeTok{key =} \DecValTok{5}\NormalTok{,}
  \AttributeTok{loudness =} \SpecialCharTok{{-}}\FloatTok{4.5}\NormalTok{,}
  \AttributeTok{mode =} \DecValTok{1}\NormalTok{,}
  \AttributeTok{speechiness =} \FloatTok{0.05}\NormalTok{,}
  \AttributeTok{acousticness =} \FloatTok{0.10}\NormalTok{,}
  \AttributeTok{instrumentalness =} \FloatTok{0.00}\NormalTok{,}
  \AttributeTok{liveness =} \FloatTok{0.15}\NormalTok{,}
  \AttributeTok{valence =} \FloatTok{0.75}\NormalTok{,}
  \AttributeTok{tempo =} \DecValTok{120}\NormalTok{,}
  \AttributeTok{time\_signature =} \DecValTok{4}\NormalTok{,}
  \AttributeTok{explicit =} \DecValTok{0}\NormalTok{,}
  \AttributeTok{track\_genre =} \StringTok{"pop"}
\NormalTok{)}

\NormalTok{nowa\_piosenka\_processed }\OtherTok{\textless{}{-}} \FunctionTok{bake}\NormalTok{(prep\_rec, }\AttributeTok{new\_data =}\NormalTok{ nowa\_piosenka)}

\NormalTok{test\_klasa }\OtherTok{\textless{}{-}} \FunctionTok{predict}\NormalTok{(rf\_fit, }\AttributeTok{new\_data =}\NormalTok{ nowa\_piosenka\_processed)}
\NormalTok{test\_prawdopodobienstwo }\OtherTok{\textless{}{-}} \FunctionTok{predict}\NormalTok{(rf\_fit, }\AttributeTok{new\_data =}\NormalTok{ nowa\_piosenka\_processed, }\AttributeTok{type =} \StringTok{"prob"}\NormalTok{)}

\FunctionTok{cat}\NormalTok{(}\StringTok{"}\SpecialCharTok{\textbackslash{}n}\StringTok{{-}{-}{-} WYNIK TESTU WDROŻENIOWEGO {-}{-}{-}}\SpecialCharTok{\textbackslash{}n}\StringTok{"}\NormalTok{)}
\end{Highlighting}
\end{Shaded}

\begin{verbatim}
## 
## --- WYNIK TESTU WDROŻENIOWEGO ---
\end{verbatim}

\begin{Shaded}
\begin{Highlighting}[]
\FunctionTok{print}\NormalTok{(test\_prawdopodobienstwo)}
\end{Highlighting}
\end{Shaded}

\begin{verbatim}
## # A tibble: 1 x 2
##   .pred_0 .pred_1
##     <dbl>   <dbl>
## 1   0.311   0.689
\end{verbatim}

\begin{Shaded}
\begin{Highlighting}[]
\FunctionTok{print}\NormalTok{(test\_klasa)}
\end{Highlighting}
\end{Shaded}

\begin{verbatim}
## # A tibble: 1 x 1
##   .pred_class
##   <fct>      
## 1 1
\end{verbatim}

\begin{Shaded}
\begin{Highlighting}[]
\NormalTok{rf\_explainer }\OtherTok{\textless{}{-}} \FunctionTok{explain\_tidymodels}\NormalTok{(}
\NormalTok{  rf\_fit,}
  \AttributeTok{data =}\NormalTok{ train\_processed }\SpecialCharTok{\%\textgreater{}\%} \FunctionTok{select}\NormalTok{(}\SpecialCharTok{{-}}\NormalTok{is\_hit),}
  \AttributeTok{y =} \FunctionTok{as.numeric}\NormalTok{(}\FunctionTok{as.character}\NormalTok{(train\_processed}\SpecialCharTok{$}\NormalTok{is\_hit)),}
  \AttributeTok{label =} \StringTok{"Random Forest"}\NormalTok{,}
  \AttributeTok{verbose =} \ConstantTok{FALSE}
\NormalTok{)}

\NormalTok{rf\_parts }\OtherTok{\textless{}{-}} \FunctionTok{predict\_parts}\NormalTok{(}\AttributeTok{explainer =}\NormalTok{ rf\_explainer,}
                          \AttributeTok{new\_observation =}\NormalTok{ nowa\_piosenka\_processed,}
                          \AttributeTok{type =} \StringTok{"break\_down"}\NormalTok{)}

\FunctionTok{plot}\NormalTok{(rf\_parts) }\SpecialCharTok{+}
  \FunctionTok{ggtitle}\NormalTok{(}\StringTok{"Wykres wodospadowy {-} Wyjaśnialność lokalna nowej piosenki"}\NormalTok{) }\SpecialCharTok{+}
  \FunctionTok{theme\_minimal}\NormalTok{()}
\end{Highlighting}
\end{Shaded}

\begin{verbatim}
## `height` was translated to `width`.
\end{verbatim}

\pandocbounded{\includegraphics[keepaspectratio]{projekt-final_files/figure-latex/unnamed-chunk-12-1.pdf}}

\textbf{Wniosek:} Test wdrożeniowy pokazuje, że model Lasu Losowego
poprawnie zakwalifikował syntetyczny profil ``piosenki pop'' jako
kandydata na hit, przypisując mu wysokie prawdopodobieństwo klasy 1.
Wykres wodospadowy (Waterfall / Break Down) dostarcza wyjaśnienia
lokalnego tej konkretnej predykcji: widać, które cechy nowej piosenki
(np. niska \texttt{instrumentalness}, wysoka \texttt{loudness}, gatunek
\texttt{pop}) ``pchnęły'' decyzję modelu w kierunku klasy pozytywnej, a
które ją hamowały. Taka interpretacja czyni model nie tylko narzędziem
predykcji, ale też źródłem rekomendacji artystycznych -- wytwórnia może
na tej podstawie identyfikować, które atrybuty nowo nagranego utworu
wymagają korekty, aby zwiększyć jego szanse komercyjne.

\subsection{Literatura}\label{literatura}

{[}1{]} Dhanaraj, R., \& Logan, B. (2005). \emph{Automatic Prediction of
Hit Songs}. Proceedings of the 6th International Symposium on Music
Information Retrieval (ISMIR 2005), 488--491.

{[}2{]} Interiano, M., Kazemi, K., Wang, L., Yang, J., Yu, Z., \&
Komarova, N. L. (2018). \emph{Musical trends and predictability of
success in contemporary songs in and out of the top charts}. Royal
Society Open Science, 5(5), 171274.
\url{https://doi.org/10.1098/rsos.171274}

{[}3{]} Chen, T., \& Guestrin, C. (2016). \emph{XGBoost: A Scalable Tree
Boosting System}. Proceedings of the 22nd ACM SIGKDD International
Conference on Knowledge Discovery and Data Mining, 785--794.
\url{https://doi.org/10.1145/2939672.2939785}

{[}4{]} Pachet, F., \& Roy, P. (2008). \emph{Hit Song Science Is Not Yet
a Science}. Proceedings of the 9th International Society for Music
Information Retrieval Conference (ISMIR 2008), 355--360.

{[}5{]} Chawla, N. V., Bowyer, K. W., Hall, L. O., \& Kegelmeyer, W. P.
(2002). \emph{SMOTE: Synthetic Minority Over-sampling Technique}.
Journal of Artificial Intelligence Research, 16, 321--357.
\url{https://doi.org/10.1613/jair.953}

\end{document}
